\documentclass[11pt]{article}

\usepackage[utf8]{inputenc}
\usepackage[T1]{fontenc}
\usepackage[margin=1in]{geometry}
\usepackage{amsmath,amssymb,amsthm}
\usepackage{booktabs}
\usepackage{graphicx}
\usepackage{microtype}
\usepackage[hidelinks]{hyperref}
\usepackage{float}

\newtheorem{lemma}{Lemma}
\newtheorem{definition}{Definition}
\newtheorem{law}{Empirical Law}

\title{\textbf{The Anatomy of a Truth Direction}\\[2pt]
\large Knowledge-Dependent Dimensionality, a Relational Law,\\
and a Convergent Category Geometry in Small Language Models}

\author{Francesco Karim Vicidomini\\
\small Independent researcher\\
\small Code: \texttt{github.com/Francesco-Marhel/TruthProbe} (one command per reported number)}

\date{July 2026\thanks{This manuscript consolidates the author's series on the
truth direction: the single-axis paper and the anatomy note (Part~I, reviewed
in an earlier form and revised following referee-style comments by B.~Nadler)
and three working notes on the cross-family mechanics and category geometry
(Part~II). Every number in both parts was regenerated under a single uniform
protocol on July 9--10, 2026. Corrections of the author's own published
interpretations are reported inside the text, as results.}}

\begin{document}
\maketitle

\begin{abstract}
B\"urger et al.\ (2024) demonstrated that truth representations in large
language models are universal across statement polarity but reside within a
multidimensional subspace. This manuscript extends that framework along
three questions: how the dimensionality of the subspace depends on the
model's knowledge, which architectural component builds the truth
direction, and what the direction is a mixture of. In Part~I, using a
training-free directional probe derived from the SVD of hidden-state
minimal pairs, we show that the dimensionality of truth is
knowledge-dependent: the signal is highly concentrated on a single axis for
behaviorally known facts (held-out AUC $0.938$) and becomes diffuse as
model knowledge decreases or material heterogeneity increases. Seven
falsification experiments sustain the one-dimensional reading on known
facts, and the same decomposition recovers the supervised polarity
direction of B\"urger et al.\ at cosine $0.959$ without polarity labels. In
Part~II, we follow the direction across four base models from two families
under pre-registered predictions. A relational empirical law emerges:
attention modules propagate truth frames they did not write, whereas
feed-forward networks oppose the frame of their current block to deposit
next-frame content; the law generalizes, and corrects, a peak-anchored
claim of our own exposed by a provenance control built against our own
instrument. Post-peak signal decay is causally attributed, on all four
models, to the SwiGLU value stream rather than its gate. Finally,
per-category truth axes form a semantically signed arrangement that
contradicts the lexical surface and converges across families
(Mantel $p = 0.0009$), indicating that the geometry of the mixture belongs
to the knowledge domain rather than to the architecture.
\end{abstract}

\section{Introduction}

Large language models produce false statements with the same fluency as
true ones, and a growing body of work shows that the truth value of a
statement is linearly readable from the residual
stream~\cite{burns2023,marks2024,li2023,zou2023}. B\"urger, Hamprecht and
Nadler~\cite{burger2024} sharpened this picture in a way that motivates
everything below: truth is \emph{universal} across statement polarity, but
it is not a single direction; probes trained on affirmative statements fail
on negated ones, and the representation lives in a two-dimensional subspace
with a dedicated polarity direction. If the truth representation has more
structure than one direction, three questions follow naturally. How much of
it fits on a single direction, and when does that reading fail? Which
architectural component builds the direction, and which degrades it? And
what is the direction made of, when the material spans many kinds of facts?

This manuscript answers the three questions in two parts, with one
instrument and one discipline. The instrument is a training-free axis: the
dominant direction of the SVD of hidden-state differences over true/false
minimal pairs, identified without labels up to one global sign
(Lemma~\ref{lem:labels}), read at cost $O(d)$ per token, and evaluated
held-out with a permutation null that prices the optimism of layer
selection. The discipline is falsification: predictions are registered
before runs, negative results are reported with the same care as positive
ones, and two corrections of our own published interpretations are narrated
in the text as results, because a claim that has survived its author's own best
attack is worth more than one that has never been attacked.

\textbf{Part~I} (one model, Qwen2.5-1.5B) establishes the object. The axis
carries a real signal on facts the model knows, in the behavioral sense
defined in Sec.~\ref{sec:data} (held-out AUC $0.938$
against a full-probe ceiling of $0.989$); the gap between axis and probe
grows as the model's behaviorally measured knowledge decreases, and grows
further when heterogeneous material forces one axis to serve many kinds of
truth: dimensionality is \emph{knowledge-dependent}. Seven falsification
experiments designed to expose a richer structure (a second dimension,
depth-global signatures, dynamics, departments, energy, phase,
three-dimensional clustering) fail to overturn the one-dimensional reading
on known facts. Blind to polarity labels, the same decomposition recovers
the supervised polarity direction $t_P$ of~\cite{burger2024} at cosine
$0.959$. An exact attention/FFN decomposition then locates the axis
anatomically: attention builds the signal in the middle layers, causally,
and the FFN degrades it after the readability peak.

\textbf{Part~II} (four base models, two families) turns the anatomy into
laws. The FFN's apparent pro-truth write at the peak and flip to anti-truth
at the next block, replicated under pre-registered predictions, is exposed
by an axis-provenance control as the special case of a \emph{relational
law}: the FFN opposes every truth frame that does not contain its own
write, while attention propagates frames it did not write
(Empirical Law~\ref{b:law:relational}); at the peak the two contributions
are near mirror images. The genuine post-peak erosion survives the
correction and is causally attributed, on all four models and with
predictions registered for the second family, to the SwiGLU value stream
rather than the gate. Two further pre-registered discriminations return
opposite verdicts, showing that scale and training recipe leave distinct
signatures: the post-peak attention--FFN tug-of-war is a trait of scale,
single-layer ignition of the signal a trait of one family. A pre-registered
negative extends cross-family: the gauge-invariant OV/QK circuit crossover
sits at a stable late relative depth on all four models and does not track
the truth peak. Finally, the axis is opened up: per-category truth axes
form a semantically signed arrangement that reorganizes with depth,
contradicts the shared lexical arrangement both families start from, and
converges across the families; the FFN writes truth along
category-dependent directions, decoding the category far above a lexical
baseline. The knowledge-dependent gap of Part~I is, in hindsight, this
mixture seen from one direction.

The two parts share every methodological choice: the same pinned dataset
revision, the same held-out protocol over pairs, the same permutation-null
estimator, the same behavioral know-rate measured before any geometry, and
the same reproducibility contract, one command per reported number.

\section{Related work}


\paragraph{Linear truth directions.} \cite{burns2023} extract a knowledge
direction without labels by enforcing logical consistency across a statement
and its negation; the direction is found by optimization. \cite{marks2024}
show truth has emergent linear structure and that difference-of-means probes
generalize as well as more complex ones. \cite{li2023} intervene along
truthful directions at inference time; \cite{zou2023} frame representation
reading top-down. Our axis shares the linear-direction hypothesis, but it is
read from the geometry of paired differences by an SVD, with no optimization;
its precise relation to supervision is stated in Lemma~\ref{lem:labels}.

\paragraph{A two-dimensional truth subspace.} Closest to us, B\"urger et
al.~\cite{burger2024} show that truth is represented not in a single direction
but in a two-dimensional subspace spanned by a general truth direction $t_G$
and a polarity-sensitive direction $t_P \sim
\mathrm{XOR}(\text{is\_true},\text{is\_negated})$, and that probes trained on
affirmative statements fail on negated ones. We confirm this structure in
Qwen2.5-1.5B and show, as a methodological addition, that $t_P$ can be
recovered \emph{without} polarity supervision.

\paragraph{Cross-layer dynamics.} Recent monitoring methods model the
evolution of hidden states across depth rather than reading a single
layer~\cite{icr2025,hsad2025}. Our depth trajectories are descriptive and
consistent with this line.

\paragraph{Beyond single-direction probes on the residual stream.} All the
methods above, ours included, assume white-box access to the residual stream,
which is not always available and is not the only usable signal. Bar-Shalom et
al.~\cite{barshalom2025} treat the full activation tensor (layers by tokens)
as an image and train a vision-transformer detector on it, across several LLMs
simultaneously, with evaluation on a broad set of hallucination benchmarks.
Methods of this kind are more expressive than any single direction and are
trained for detection; our axis sits at the opposite end of the same
cost/expressivity curve, as the cheapest structural point, and the comparison
we draw is against that curve, not against any single method.

\paragraph{Superposition and dictionaries.} Superposition~\cite{elhage2022}
motivates decomposing the residual stream into overcomplete
dictionaries~\cite{cunningham2023,bricken2023}. This is the operation we
explicitly do \emph{not} perform. The bottom line, argued in
Secs.~\ref{sec:discussion} and~\ref{sec:why}, is that the two approaches
answer different questions and neither subsumes the other: a dictionary
answers ``which features are active'' at cost $O(d\,m)$ and after dedicated
training; the axis answers ``how much truth a single direction carries'' at
cost $O(d)$ with no training. When feature identity matters, dictionaries are
the right tool; for a cheap runtime signal on facts the model knows, the axis
suffices. We make no claim of superior accuracy over dictionary methods.

\paragraph{MLPs as key-value writers.} The memory-editing
literature~\cite{meng2022} attributes to transformer MLPs the physiology of
key-value stores whose writes deposit factual content into the residual
stream. The relational law of Sec.~\ref{b:sec:law} is a dynamical face of that
picture: the component that deposits the next frame's material is, per these
measurements, the one that opposes the current frame.

\paragraph{Rotational dynamics of factual selection.} A recent
preprint~\cite{rotational2026} reports that factual constraint processing
separates alternatives by rotation, with active suppression of the correct
continuation at intermediate depth. The propagate/overwrite asymmetry
supplies a candidate microscopic writer for that macroscopic observation.

\paragraph{Attention-head specialization.} The head-taxonomy line of
work~\cite{headgenome} motivated the pre-registered test of whether static,
gauge-invariant circuit norms predict where truth lives (they do not,
Sec.~\ref{b:sec:negatives}) and defines the per-head follow-up of
Sec.~\ref{b:sec:leads}.

\paragraph{Cross-model representation convergence.} The convergence of the
category arrangement across families (Sec.~\ref{b:sec:mixture}) is a signed,
truth-specific instance of the broader phenomenon that different models
trained on the same world converge to similar
representations~\cite{platonic2024}.


\part*{Part I: The Axis on One Model}
\addcontentsline{toc}{part}{Part I}
\section{Scope and evidence}
\label{sec:scope}

Everything in this paper concerns one small base model, Qwen2.5-1.5B, and
sentence-level factual statements. Results come in two evidential tiers, and
we keep them distinct throughout. \textbf{Tier 1} (Secs.~\ref{sec:results}
and~\ref{sec:robust}): the positive findings and the falsification battery,
evaluated held-out with permutation nulls. \textbf{Tier 2}
(Sec.~\ref{sec:anatomy}): the anatomy, a mechanistic follow-up with the same
protocol but smaller margins in places, single datasets for some entries, and
no falsification battery of its own; where a statement there is
interpretation rather than measurement, it is marked as such. Cross-model
generalization is not claimed here; what is known about it from subsequent
work is summarized, with numbers, in Part~II.

\section{Method}
\label{sec:method}

\subsection{Extraction}
\label{sec:extraction}

Informally: a transformer reads a sentence and, at every layer, updates an
internal vector summary of the text so far; we record that summary at the
position of the last token, where the sentence is complete, at every layer,
and study the geometry of these vectors.

\paragraph{Pooling and the token-identity confound.} All states are read
at the last token, and in minimal pairs the flipped content word sits at
or immediately before that position: token identity is therefore a
candidate confound for everything downstream. Three measurements bound
it. First, the per-layer curves start at chance in every scan (layer-$0$
AUC $= 0.500$ on all four models): a token-identity axis would separate
already at the embedding, and it does not; the signal emerges only with
depth, that is, with context integration. Second, all AUCs are held-out
across pairs whose targets are disjoint, which a pure target-identity
direction would not survive at the measured levels. Third, the category
analysis carries its own lexical control: the block-$2$ baseline decodes
the relation at $39$--$42\%$ against $74$--$82\%$ for the peak-adjacent
FFN writes. The extraction tool also supports mean pooling
(\texttt{--pool mean}); the last-token results are the canonical ones,
and the mean-pooling replication is listed as a cheap robustness run
rather than assumed.


Formally, we load Qwen2.5-1.5B frozen. The model has $28$ transformer blocks,
hence $29$ hidden-state levels (the embedding plus each block output), each a
vector of width $d = 1536$, the width of the residual stream. Per sentence we
run one forward pass and take the last-token state at every level. For $N$
sentences this yields $H \in \mathbb{R}^{N \times 29 \times d}$. Once an axis
is fixed, scoring a new state is a single dot product, hence the reading cost
of $O(d)$ per token quoted throughout.

\subsection{Data: minimal pairs}
\label{sec:data}

A \emph{minimal pair} is a pair of sentences with the same subject and the
same template that differ in exactly one content word, such that one sentence
is true and the other false: ``The Sun is a star.''\ vs.\ ``The Sun is a
planet.''\ Both sentences are affirmative. Within a pair, the difference of
hidden states cancels the shared topic and surface form and leaves the truth
signal.

Two sources are used. The \emph{curated pairs} are a built-in set of $45$
pairs of common facts (capitals, chemical symbols, authorship, currencies)
that the model demonstrably reproduces in open-ended completion
(Sec.~\ref{sec:future}); they play the role of a known-facts control:
behaviorally verified and token-matched, hence also clean, in the sense of
free of the sentence-format confounds discussed in Sec.~\ref{sec:limits}.
The external datasets are CounterFact~\cite{counterfact} and
TruthfulQA~\cite{truthfulqa} ($250$ pairs, pinned revisions in
Sec.~\ref{sec:repro}), which contain many relations a $1.5$B model plausibly
does not know.

To quantify this, we define the \emph{know-rate}: the fraction of facts the
model reproduces correctly in open-ended greedy completion. Over $200$
CounterFact prompts, Qwen2.5-1.5B has a know-rate of $30.0\%$ ($61/200$),
which grounds the ``plausibly does not know'' reading; the know-rates of the
other models appear in Part~II.

Precisely: a prompt counts as known when the greedy completion
(\texttt{do\_sample=False}, at most $8$ new tokens) either begins with the
true target or contains it preceded by a space, case-insensitively
(\texttt{behav\_check.py}). Three properties of this estimator matter for
what follows. It is a \emph{lower bound}: a model may know a fact and
still continue differently within eight tokens. It is \emph{greedy}: no
sampling, so the estimate is deterministic given the seed that selects the
prompts. And it inherits tokenizer and completion-style differences
between families, so cross-family comparisons carry a systematic bias
component, while within a family the bias is shared and cancels. This is
why the dimensionality relation is stated as within-family monotone with a
cross-family offset, and why the ordering Llama-3.2-1B ($35.5\%$) above
Qwen2.5-1.5B ($30.0\%$) is neither exotic nor load-bearing: parameter
count is not knowledge (the two models differ in training data, tokenizer,
and completion style), and no claim in this manuscript rests on comparing
know-rates across families.



\paragraph{Why a flipped predicate and not a negation.} The false counterpart
is built by flipping a content word (``planet'' for ``star''), not by
inserting ``not''. The choice is deliberate, for three reasons. First,
negation is not a neutral edit: it activates a dedicated polarity direction in
the representation~\cite{burger2024}, so a negated counterpart would entangle
the truth signal with the polarity signal. Our design keeps both sentences
affirmative and isolates truth; polarity is then studied separately and
deliberately in the $2{\times}2$ design of Sec.~\ref{sec:poldesign}. Second,
inserting ``not'' changes the token count and shifts all subsequent positions,
weakening the token-level matching that makes the intra-pair difference
cancel shared structure. Third, the two designs are complementary rather than
competing: the negated forms are exactly the NT/NF cells of the polarity
design, so the alternative is not ignored but treated as its own object of
study.

\subsection{The axis, and where labels enter}
\label{sec:axis}

We form the intra-pair difference matrix $D$ with rows
$h_{\text{true}} - h_{\text{false}}$ at a given layer, and take its SVD,
\begin{equation}
D = U S V^{\top}, \qquad v_{1} = V_{1},
\end{equation}
where $v_1$, the top right-singular vector, is the dominant direction of the
differences. The role of supervision in this construction can be stated
exactly.

\begin{lemma}[Label-free identification of the axis, up to sign]
\label{lem:labels}
Let each pair $i$ contribute a row $d_i^{\top}$ to $D$, where
$d_i = s_i\,(h_i^{A} - h_i^{B}) \in \mathbb{R}^{d}$ is a column vector,
$h_i^{A}, h_i^{B}$ are the two hidden states of the pair, and
$s_i \in \{\pm 1\}$ encodes which sentence of the pair is designated true.
Then $D^{\top} D = \sum_i d_i d_i^{\top}$ does not depend on $(s_i)$; hence
the right singular vectors of $D$, and in particular the span of $v_1$, are
invariant to the within-pair labels.
\end{lemma}

One implicit choice is declared here rather than discovered by a referee:
the SVD is applied to the raw difference rows, so $v_1$ is a
\emph{norm-weighted} principal direction, with pair $i$ contributing
weight proportional to $\lVert d_i \rVert^{2}$. If difference norms
correlate with sentence length or with the state norm at the reading
position, long or high-norm pairs dominate the fit. The control was
executed on all four models at their peaks
(\texttt{axis\_norm\_check.py}): the row-normalized refit agrees with the
raw axis at $|\cos| = 0.984$, $0.980$, $0.974$, $0.966$ (Qwen2.5-1.5B,
Qwen2.5-3B, Llama-3.2-1B, Llama-3.2-3B), the held-out AUC changes by at
most $+0.016$, and is in fact slightly higher under normalization on every
model, so the axis does not ride on high-norm pairs; difference norms
correlate only weakly, and negatively, with sentence length ($r$ between
$-0.12$ and $-0.23$). The axis is not a norm-weighting artifact.


\begin{proof}
$d_i d_i^{\top} = (-d_i)(-d_i)^{\top}$ for every $i$, so each summand of the
Gram matrix, and therefore the Gram matrix and its eigenvectors, is unchanged
under any sign assignment.
\end{proof}

\paragraph{What label-free does and does not mean.} Three levels of
supervision must be kept apart. (i) \emph{Dataset construction} is
supervised: writing a minimal pair requires knowing which completion is
factually correct, and no method that starts from such pairs can call
itself unsupervised at this level. (ii) \emph{Axis estimation} consumes
only the contrast structure: by Lemma~\ref{lem:labels}, the span of $v_1$
is invariant to which element of each pair is designated true, so an
adversary who scrambled every within-pair label would hand the algorithm
the same axis. The information the estimator uses is ``these two
statements differ in truth value'', not ``this one is true''.
(iii) \emph{Orientation and evaluation} are supervised: one global sign
bit, and held-out labels for scoring. Accordingly, this manuscript says
\emph{training-free} of the pipeline, \emph{label-free up to sign, given
the paired design} of the estimator, and reserves \emph{unsupervised} for
the polarity recovery of Part~I, where the recovered factor was never
encoded in any label handed to the algorithm.


Labels enter the pipeline at exactly two points. First, a single global
orientation bit (``the true side is positive''), fit on training folds,
fixes the sign of $v_1$. Second, the evaluation itself: AUC is defined with
respect to labels. Everything upstream of these two points, the pairing, the
differences, and the axis, uses no labels. We therefore describe the axis as
\emph{training-free} (no optimization of any kind) and \emph{label-free up to
one sign bit} (Lemma~\ref{lem:labels}), and we reserve the word
``unsupervised'' for the polarity recovery of Sec.~\ref{sec:polarity}, where
the recovered structure (polarity) is never labeled at all. The lemma also
explains, in advance, the behavior of the permutation null below: permuting
which sentence of a pair is called true leaves the span of the axis
unchanged and destroys only the orientation and the test labels, which is why
the null concentrates near chance.

\subsection{A falsifiable evaluation protocol}
\label{sec:protocol}

The central methodological point is that the axis, its calibration, the
orientation bit, \emph{and} the choice of peak layer (the layer whose held-out AUC is maximal, chosen on training folds) must be fit on training
data only. We use $k$-fold cross-validation \emph{over pairs}: a pair is
never split across folds, which would leak the shared topic. Everything is
fit on the training folds and we report held-out AUC averaged across folds.
Pairs from the same CounterFact relation template may still fall in
different folds; generalization across topics is therefore not certified by
this protocol, and we return to the point in Sec.~\ref{sec:limits}.

Part of the apparent accuracy of any layer-scanned pipeline is selection: the
best of $29$ layers looks good by chance. The \emph{permutation null}
quantifies exactly that contribution. Let $A_{\mathrm{obs}}$ be the observed
maximum held-out AUC over layers (AUC: the area under the ROC curve; equivalently, the probability that a
randomly chosen true statement scores above a randomly chosen false one,
ties counted as one half). For $b = 1, \dots, B$ we randomly swap
true/false \emph{within} each pair (the labels become noise while the pair
structure is preserved), rerun the identical pipeline including the layer
scan, and record the maximum held-out AUC $A_b$. The estimated $p$-value is
the standard conservative estimator
\begin{equation}
\hat{p} \;=\; \frac{1 + \#\{\,b : A_b \ge A_{\mathrm{obs}}\,\}}{B + 1}.
\end{equation}
With $B = 200$ and no null value reaching the observed statistic,
$\hat{p} = 1/201 \approx 0.005$, the value reported in
Table~\ref{tab:main}. As a reference point we also fit an
$\ell_2$-regularized logistic probe on the raw $d$-dimensional states, held
out over the same pair folds.

\subsection{A complex-plane extension}
\label{sec:complexmethod}

We additionally form a second axis $v_2$, the SVD residual orthogonal to
$v_1$, treat the projection as $z = c_1 + i\,c_2$ with calibrated real and
imaginary parts, and test whether magnitude- or phase-based scores improve
separation over the real axis alone. They do not (Sec.~\ref{sec:phase}); the
negative finding is kept because it is itself informative.

\subsection{Polarity design}
\label{sec:poldesign}

To test whether a genuine second dimension exists, we use a $2{\times}2$
design: each fact in four forms, affirmative-true (AT), affirmative-false
(AF), negated-true (NT), negated-false (NF), giving the XOR structure
of~\cite{burger2024}. We compute the affirmative and negated truth directions
$t_A = \overline{AT} - \overline{AF}$, $t_N = \overline{NT} - \overline{NF}$,
the general direction $t_G \propto t_A + t_N$, and the supervised polarity
direction $t_P \propto t_A - t_N$. For the recovery test, we run the
difference SVD on minimal pairs that \emph{mix} the two polarities, so the
SVD never sees the polarity label, and compare its leading directions to the
supervised $t_P$.

\subsection{Residual decomposition and its correctness gate}
\label{sec:decomp}

Within a block the residual stream is updated additively,
\begin{equation}
h_{\ell+1} \;=\; h_{\ell} + a_{\ell} + f_{\ell},
\end{equation}
where $a_{\ell}$ is the attention contribution and $f_{\ell}$ the FFN
contribution, both captured by forward hooks. The residual stream studied
everywhere else in this paper is the sum; the anatomy of
Sec.~\ref{sec:anatomy} reads the two addends separately.

Every anatomy tool embeds an \textbf{identity check}: it verifies that
$a_{\ell} + f_{\ell}$ reconstructs the residual delta $h_{\ell+1} - h_{\ell}$.
The verdict uses the \emph{median} relative error across (sentence, layer)
rather than the maximum: where a block adds almost nothing the delta is near
zero and the ratio explodes, so the maximum is fragile while the median
reports the typical case. Observed median error: about $1.0 \times 10^{-7}$
on both datasets, so the decomposition is exact and the anatomy numbers are
trustworthy. (The single inflated layer, block $27$, has a near-zero residual
delta and is excluded from interpretation.) Anatomy runs use
\texttt{float32}: the identity fails spuriously in \texttt{bfloat16} because
subtracting large residual states to form small per-block deltas suffers
catastrophic cancellation.

\section{Results I: signal, dimensionality, polarity}
\label{sec:results}

All numbers are on Qwen2.5-1.5B with last-token pooling; AUC is out-of-fold
unless stated.

\subsection{A real but bounded truth signal}
\label{sec:signal}

Table~\ref{tab:main} reports the held-out evaluation at the peak layer (layer
$16$). On the curated pairs the training-free axis reaches AUC $0.938$ with
permutation $\hat{p} = 0.005$ (null mean $0.588$, $95$th percentile $0.664$):
the geometric axis aligns with truth well beyond what layer selection alone
can manufacture. The signal is therefore real. It is also \emph{bounded}: a
full linear probe on the same raw states reaches $0.989$. The probe sees
information the single axis discards.

\begin{table}[t]
\centering
\caption{Held-out AUC ($5$-fold CV over pairs) at layer $16$, Qwen2.5-1.5B.
``1D'' is the real axis; ``MAG'' the magnitude-gated polar score; ``PHASE''
the angular-deviation score; ``2D'' a logistic on
$(\mathrm{Re},\mathrm{Im})$; ``Probe'' an $\ell_2$ logistic on raw states.
The last column is the label-permutation $\hat{p}$ for the 1D axis.}
\label{tab:main}
\begin{tabular}{lccccccc}
\toprule
Dataset & pairs & 1D & MAG & PHASE & 2D & Probe & $\hat{p}$ \\
\midrule
Curated pairs                    & $45$  & $0.938$ & $0.936$ & $0.928$ & $0.931$ & $0.989$ & $0.005$ \\
CounterFact only                 & $250$ & $0.792$ & $0.790$ & $0.771$ & $0.795$ & $0.854$ & $0.005$ \\
CounterFact / TruthfulQA mix     & $250$ & $0.638$ & $0.640$ & $0.676$ & $0.651$ & $0.841$ & $0.005$ \\
\bottomrule
\end{tabular}
\end{table}

\subsection{The gap scales with model knowledge, and with heterogeneity}
\label{sec:dim}

The rows of Table~\ref{tab:main} differ not only in difficulty but in the
\emph{size of the probe gap}, and the three rows separate two distinct
effects.

The first effect is knowledge. On curated facts the model demonstrably knows,
the axis ($0.938$) trails the full probe ($0.989$) by $0.051$. On CounterFact
alone, where the measured know-rate is $30.0\%$ (Sec.~\ref{sec:data}) and
many relations are obscure (for example, the mother tongue of a minor
historical figure), the separability of the axis falls to $0.792$ and the gap
widens to $0.062$ (probe $0.854$). Less knowledge, lower separability, and a
modestly larger gap.

The second effect is heterogeneity. On the CounterFact/TruthfulQA mix the gap
jumps to $0.203$ (axis $0.638$, probe $0.841$), far more than either
ingredient alone would suggest. A single direction fit across heterogeneous
material must serve all of it at once, while the full probe is free to
combine directions; the excess gap on the mix therefore measures how much the
truth direction itself varies with the material. Subsequent work confirms
this reading independently Part~II: the
per-category truth axes form a structured, semantically signed geometry, and
forcing one axis onto a mixture of categories is exactly what the mix row
pays for.

We read the knowledge effect as a statement about \emph{dimensionality}. When
the model knows a fact, its truth content is concentrated and a single
direction captures most of it. When the model does not know a fact, the
truth-relevant signal is weaker and more diffuse, and one direction recovers
less of it. The dimensionality of the truth representation is thus not fixed
but \emph{knowledge-dependent}; the mix row adds the complementary statement
that it is also \emph{category-dependent}. This joint framing is, to our
knowledge, stated in these terms by neither the linear-probe nor the
dictionary-learning literature, and it is the contribution we consider most
novel.

\subsection{The complex plane does not help; ``phase'' is the sign flip}
\label{sec:phase}

On single-polarity data the second dimension adds nothing: at layer $16$ the
2D logistic ($0.931$) does not exceed the 1D axis ($0.938$) on curated pairs,
and on CounterFact it gains $+0.003$ ($0.795$ vs.\ $0.792$). This is
structural. Minimal pairs cancel the topic, so the differences concentrate
along one direction; the SVD places the entire signal in $v_1$, and $v_2$ is
residual variance, noise rather than a second truth factor. A controlled
synthetic check confirms that the apparent ``rotation'' of falsehoods is an
artifact: with a pure sign flip and zero genuine rotation, an angular (phase)
score still separates classes, because a state at $\mathrm{Re} < 0,\
\mathrm{Im} \approx 0$ has phase $\approx \pi$ by definition. The earlier
intuition of truth-as-phase has no operational content here: it reduces to
the sign of the real axis re-expressed in polar coordinates. We retract that
framing.

\subsection{Recovery of the polarity direction without polarity labels}
\label{sec:polarity}

A genuine second dimension does exist, but only when the data span both
polarities, and it is the polarity direction of~\cite{burger2024}, not a
phase. Table~\ref{tab:polarity} shows the diagnostic: an axis fit on
affirmative statements classifies affirmative held-out facts at $0.981$ but
\emph{anti-classifies} negated ones at $0.079$, a near-total flip, while the
general direction $t_G$ recovers the negated facts at $0.954$. The truth
direction genuinely rotates with polarity ($\cos(t_A, t_N) = -0.415$).

The methodological result is the recovery. When we run the difference SVD on
\emph{mixed-polarity} minimal pairs, blind to the negation label, its leading
direction aligns with the supervised polarity direction $t_P$ at cosine
$0.959$ (layer $15$), and $t_P$ lands in the first principal component, while
alignment with $t_G$ is only $0.157$. In other words, when both polarities
are present, the dominant axis of the differences \emph{is} polarity, and the
SVD finds it without supervision; by Lemma~\ref{lem:labels} the recovery
cannot have used the labels even implicitly. Where~\cite{burger2024}
constructs $t_P$ by separating affirmative from negated statements by hand,
the same direction emerges here from geometry alone.

\begin{table}[t]
\centering
\caption{Polarity analysis (leave-facts-out CV), layer $14$ chosen on
affirmative data; and unsupervised $t_P$ recovery at layer $15$. AUC below
$0.5$ indicates a sign flip.}
\label{tab:polarity}
\begin{tabular}{lc}
\toprule
Quantity & Value \\
\midrule
$\cos(t_A, t_N)$ & $-0.415$ \\
Affirmative axis on affirmative (held-out AUC) & $0.981$ \\
Affirmative axis on \emph{negated} (held-out AUC) & $0.079$ \\
General direction $t_G$ on negated (held-out AUC) & $0.954$ \\
\midrule
SVD vs.\ supervised $t_P$: $\max_k |\cos(\mathrm{SVD}_k, t_P)|$ & $0.959$ \\
\quad (principal component carrying $t_P$) & PC\,$1$ \\
$\cos(\mathrm{SVD}_1, t_G)$ & $0.157$ \\
\bottomrule
\end{tabular}
\end{table}


\section{Robustness: seven falsification experiments}
\label{sec:robust}

The first half of this paper established a positive claim. We now stress-test
it. The claim under test is:

\begin{quote}
\emph{In Qwen2.5-1.5B, the truth value of a statement is carried by a single
(one-dimensional) direction, the position axis obtained from the SVD of
paired true/false differences, on facts the model knows; the second and
higher dimensions add no separation.}
\end{quote}

A claim of one-dimensionality is falsifiable in a specific way: exhibit a
second, third, or depth-distributed coordinate that separates true from false
better than the single axis. We construct seven such coordinates and test
each.

\paragraph{Two kinds of test, kept distinct.} We separate \emph{controlled
experiments} from \emph{visual observations}, and label every result as one
or the other. Controlled experiments (experiments 1--4) produce a held-out
AUC under $k$-fold cross-validation over pairs and a label-permutation null
that quantifies selection bias; a result is meaningful only if it clears that
null. Visual observations (experiments 5--7) do not classify; they project
the states and measure a single descriptive statistic, the per-coordinate
separation AUC or the centroid-separation ratio, so the conclusion follows
from a number attached to a picture, not from a fitted model. Throughout, the
model is frozen Qwen2.5-1.5B, last-token pooling, peak layer $16$ unless
stated.

Table~\ref{tab:seven} summarizes all seven with their real numbers; the prose
below gives the reasoning for each.

\begin{table}[t]
\centering
\caption{Seven experiments attempting to falsify the one-dimensional truth
axis. ``Exp.'' = controlled experiment (held-out CV + permutation null);
``Obs.'' = visual observation (descriptive separation statistic). None
overturns the claim.}
\label{tab:seven}
\small
\begin{tabular}{p{0.5cm}p{4.6cm}p{1.1cm}p{4.6cm}p{1.4cm}}
\toprule
\# & Richer structure hypothesized & Type & Result (real numbers) & Verdict \\
\midrule
1 & A complex/phase second dimension (truth as rotation) & Exp. & PHASE AUC $0.928$ vs.\ 1D $0.938$ (curated); reduces to the real-axis sign flip & falsified \\
2 & Truth is global across depth (fixed-axis signature) & Exp. & global $0.923$ vs.\ local $0.938$ ($-0.015$), $\hat{p} = 0.0099$; persistence gap $+0.04$ & falsified \\
3 & Truth is dynamic, a ``butterfly effect'' in the unsaturated shade, on hard data & Exp. & global $0.642$ vs.\ local $0.627$ ($+0.015$), $\hat{p} = 0.0099$; commitment true $2.83 \to 5.91$ vs.\ false $2.65 \to 5.86$ & falsified \\
4 & Truth is distributed across layer ``departments'' & Exp. & diagnostic curve is a smooth ramp (no steps); global $0.646$ vs.\ local $0.627$ ($+0.019$), $\hat{p} = 0.0099$ & falsified \\
5 & The argument $\theta$ (phase) carries truth & Obs. & $\theta$ AUC $0.504$ / $0.584$ / $0.556$ (CF / TQA / mix) & falsified \\
6 & The energy $|z|$ (commitment) carries truth & Obs. & $|z|$ AUC $0.560$ / $0.558$ / $0.543$ & falsified \\
7 & True/false form distinct 3D clumps & Obs. & centroid ratio $0.55$ / $0.51$ / $0.48$ ($<1$, overlapping); $\mathrm{Re}$ alone $0.795$ / $0.656$ & falsified \\
\bottomrule
\end{tabular}
\end{table}

\paragraph{1. The complex plane.} The original framing promoted the second
SVD direction to an imaginary axis and read truth as a phase rotation. If
real, a phase- or magnitude-based score should beat the real axis. It does
not: on curated pairs the phase score reaches AUC $0.928$ against $0.938$ for
the real axis, and a controlled synthetic check shows that with a pure sign
flip and \emph{zero} genuine rotation an angular score still separates,
because a state at $\mathrm{Re} < 0$ has phase $\approx \pi$ by definition.
The ``rotation'' is the sign of the real axis re-expressed in polar
coordinates. The complex framing is decorative; we retract it. (This is the
negative result of Sec.~\ref{sec:phase}, restated as the first falsification
experiment.)

\textbf{2. A global depth signature.} If truth flows through the network, the full trajectory $[\Delta_0, \dots, \Delta_{28}]$ should beat the single best layer, where the signed deviation at each layer $\ell$ is defined as $\Delta_\ell = b_\ell - \frac{1}{2}$, with $b_\ell = \sigma(\operatorname{Re}(z_\ell))$ representing the calibrated projection passed through a sigmoid along a fixed axis. On curated pairs it does not: the depth-signature classifier reaches held-out AUC 0.923 against 0.938 for the single layer (-0.015), and although the signature is real (permutation $\hat{p} = 0.0099$), it carries no advantage. The crossing-persistence gap between false and true is $+0.04$, flat (where for each sentence, the first layer at which the trajectory crosses below zero and the fraction of subsequent layers it stays below; the gap is the difference of this persistence between false and true statements). No propagation.


\paragraph{3. A dynamic butterfly effect in the shade.} The signed deviation
saturates ($\sigma$ compresses large values to $\pm\tfrac12$), so a
propagating \emph{uncertainty} could be invisible to it. We re-ran on hard
data (CounterFact/TruthfulQA mix) with the unsaturated features
$[\mathrm{Re}, |z|]$, the ``shade''. The depth signature now beats the single
layer, but by $+0.015$ ($0.642$ vs.\ $0.627$, $\hat{p} = 0.0099$): real and
marginal. The decisive descriptive measure kills the butterfly effect: the
commitment magnitude grows identically for both classes (true
$2.83 \to 5.91$, false $2.65 \to 5.86$); false statements do not hesitate
more than true ones, and there is no propagating divergence.

\paragraph{4. Layer departments.} If different layer bands perform different
jobs (surface, semantics, logic), a single semantic axis re-read everywhere
wastes most layers; giving each layer its \emph{own} native probe should
recover a distributed signal, and the per-layer separation curve should show
steps. It shows a smooth ramp instead: the per-layer native separation rises
continuously to a single peak at layer $16$, with no plateaus and no jumps.
The per-layer signature beats the single layer by $+0.019$ ($0.646$ vs.\
$0.627$, $\hat{p} = 0.0099$), again real and marginal. No departments; the
native-axis ``cure'' improves the fixed axis but does not change the ceiling.

\paragraph{5--6. Energy and argument.} Decomposing each state into position
($\mathrm{Re}$), energy ($|z|$), and argument ($\theta$) and measuring how
well each separates true from false (Fig.~\ref{fig:coords}): only position
carries signal. The energy separates at $|z|$ AUC $0.560$/$0.558$/$0.543$
(CounterFact/TruthfulQA/mix): true and false are equally committed. The
argument separates at $\theta$ AUC $0.504$/$0.584$/$0.556$, essentially
chance; what little it shows is the echo of the $\mathrm{Re}$ sign, not
independent information. The phase, the very coordinate the complex framing
made central, is empty.

\paragraph{7. Three-dimensional clumping.} If truth were a richer object,
true and false should form distinct clusters in the three-dimensional
$(v_1, v_2, v_3)$ space. They do not (Fig.~\ref{fig:plane}). The
centroid-separation ratio, the distance between the true and false centroids
divided by the average cloud spread, is $0.55$/$0.51$/$0.48$, all below $1$:
the clouds overlap. Crucially this low ratio is \emph{not} evidence against
the axis; it is the one-dimensional signal seen through two noise dimensions.
Projected onto position alone, the same data separate at AUC $0.795$
(CounterFact). Adding the orthogonal dimensions dilutes, it does not enrich.

\subsection{What survives}
\label{sec:survive}

Three statements resist all seven experiments.

\emph{Truth is one-dimensional and lives on the position axis.} Across every
decomposition, only $\mathrm{Re}$ separates; energy, phase, the second and
third SVD directions, and the full-depth trajectory add at most $+0.02$ AUC
and usually nothing.

\emph{Separability scales with knowledge.} On CounterFact, where the model
often knows the fact, the position axis separates at AUC $0.795$; on the more
ambiguous TruthfulQA at $0.656$; on the diffuse mix the centroids nearly
coincide. The clearer the model's knowledge, the more one-dimensional and
separable the truth representation.

\emph{On unknown facts, the classes show low separability.} Where the model
does not know the fact, true and false representations do not form distinct
clusters but overlap in space, intuitively a dense, uniform distribution
rather than two distinct clumps. We state this as \emph{low class
separability / distribution overlap} (centroid ratio $<1$), the measured
counterpart of the knowledge-dependent dimensionality of Sec.~\ref{sec:dim}.

\begin{figure}[t]
\centering
% Upload counterfact_coords.png and truthfulqa_coords.png alongside the .tex.
\begin{minipage}{0.49\textwidth}\centering
\includegraphics[width=\textwidth]{counterfact_coords.png}\\[2pt]
{\footnotesize (a) CounterFact: $\mathrm{Re}$ $0.795$, $|z|$ $0.560$, $\theta$ $0.504$}
\end{minipage}\hfill
\begin{minipage}{0.49\textwidth}\centering
\includegraphics[width=\textwidth]{truthfulqa_coords.png}\\[2pt]
{\footnotesize (b) TruthfulQA: $\mathrm{Re}$ $0.656$, $|z|$ $0.558$, $\theta$ $0.584$}
\end{minipage}
\caption{Experiments 5--6 and knowledge scaling, layer $16$. Per-coordinate
separation AUC. On both datasets only the position axis ($\mathrm{Re}$)
separates true from false; energy ($|z|$) and argument ($\theta$) are at
chance. The position separation falls from $0.795$ (CounterFact, often-known
facts) to $0.656$ (TruthfulQA, more ambiguous), while the one-dimensional
structure is unchanged.}
\label{fig:coords}
\end{figure}

\begin{figure}[t]
\centering
% Upload counterfact_plane.png and counterfact_3d.png alongside the .tex.
\begin{minipage}{0.49\textwidth}\centering
\includegraphics[width=\textwidth]{counterfact_plane.png}
\end{minipage}\hfill
\begin{minipage}{0.49\textwidth}\centering
\includegraphics[width=\textwidth]{counterfact_3d.png}
\end{minipage}
\caption{Experiment 7, CounterFact at layer $16$. Left: the 2D plane; the
true (blue) and false (red) centroids (diamonds) are close relative to the
cloud spread (centroid-separation ratio $0.55$, overlapping). Right: the
three-dimensional $(v_1, v_2, v_3)$ scatter shows a gradient along $v_1$
(position) but no distinct clumps. The separation is confined to the
horizontal position axis; the extra dimensions are noise.}
\label{fig:plane}
\end{figure}

\begin{figure}[t]
\centering
% Upload truthfulqa_plane.png and truthfulqa_3d.png alongside the .tex.
\begin{minipage}{0.49\textwidth}\centering
\includegraphics[width=\textwidth]{truthfulqa_plane.png}
\end{minipage}\hfill
\begin{minipage}{0.49\textwidth}\centering
\includegraphics[width=\textwidth]{truthfulqa_3d.png}
\end{minipage}
\caption{Experiment 7, TruthfulQA at layer $16$. The same picture on more
ambiguous facts: centroid-separation ratio $0.51$ (overlapping), no
three-dimensional clumps. Lower knowledge, lower separability, identical
one-dimensional structure.}
\label{fig:plane_tqa}
\end{figure}


\section{Results II: the anatomy of the axis}
\label{sec:anatomy}

The falsification battery established \emph{that} truth is one-dimensional on
known facts. This section asks \emph{how the residual stream comes to carry
it}: which architectural component, attention or the FFN, produces, holds, or
degrades the signal, layer by layer. The results here are Tier 2 in the sense
of Sec.~\ref{sec:scope}: same protocol and safeguards, but smaller margins in
places, and interpretations are marked as such. The decomposition and its
correctness gate are defined in Sec.~\ref{sec:decomp}.

\subsection{Per-layer readability: a stable alternation}
\label{sec:alternation}

For each layer we fit the truth axis on $a_{\ell}$ and on $f_{\ell}$
separately and measure held-out AUC, repeating over $5$ cross-validation
seeds. A per-layer difference $\mathrm{AUC}(a_\ell) - \mathrm{AUC}(f_\ell)$
is counted as real only if its magnitude exceeds both $2\sigma$ (its own
across-seed standard deviation) and a floor of $0.03$. Table~\ref{tab:peak}
shows the peak block.

\begin{table}[h]
\centering
\caption{Truth AUC by component at the readability peak (block $15$),
CounterFact, mean $\pm$ std over $5$ seeds. Residual = the summed stream
(reference).}
\label{tab:peak}
\begin{tabular}{lccc}
\toprule
 & residual & attention & FFN \\
\midrule
AUC & $0.792 \pm 0.003$ & $0.845 \pm 0.003$ & $0.734 \pm 0.004$ \\
\bottomrule
\end{tabular}
\end{table}

The alternation survives re-seeding and reproduces structurally across both
datasets. On CounterFact, attention stably dominates readability at layers
$\{6,7,10,12,13,14,15,17,18\}$ and the FFN at
$\{8,11,16,19,20,21,22,23,25,26\}$. At the peak (block $15$), where the truth
axis is read, attention leads by $+0.110 \pm 0.004$. \textbf{Interpretation
(not yet causal):} attention appears to carry the truth signal most strongly
through the middle band, including the peak. Readability, however, is not
necessity, which Sec.~\ref{sec:ablation} tests directly, and which changes
the picture.

\subsection{The expanded FFN space does not help}
\label{sec:expanded}

The FFN expands its input to $8960$ dimensions, applies SwiGLU, and
re-compresses it. One might expect the truth signal to be more cleanly
readable inside that larger, putatively sparser space, before compression. We
capture the expanded activation (the input to \texttt{down\_proj}) and fit
the truth axis there. In this $p \gg n$ regime the permutation null is
elevated, so only the margin over the null counts: at the peak the expanded
axis clears its null cleanly (curated $0.953$ vs.\ null $95$th percentile
$0.669$, margin $+0.284$, $\hat{p} = 0.0099$; CounterFact $0.816$ vs.\
$0.561$, margin $+0.254$, $\hat{p} = 0.0099$). \textbf{But relative to the
compressed residual at the same layer, the expanded space gains essentially
nothing}: $+0.017$ (curated) and $+0.023$ (CounterFact). Raw architectural
expansion does not undo superposition on these facts. The implication for
future work is concrete: sparsity must be \emph{imposed}, for example by a
trained sparse autoencoder, not merely afforded by a wider space.

\subsection{Causal ablation: attention builds, the FFN erodes}
\label{sec:ablation}

To separate readability from necessity we zero a component's contribution
over a band of layers (a hook returns zeros, so $h + 0 = h$), let the effect
propagate, and measure the collapse of truth AUC at a downstream readout
(Table~\ref{tab:abl}).

\begin{table}[h]
\centering
\caption{Causal ablation on CounterFact ($250$ pairs). Drop $=$ intact $-$
ablated; positive drop $=$ the component is causally necessary.}
\label{tab:abl}
\begin{tabular}{llccc}
\toprule
band (layers) & readout & intact & attention-off (drop) & FFN-off (drop) \\
\midrule
$11$--$15$ (middle)     & block $15$ & $0.793$ & $0.487\;(+0.306)$ & $0.819\;(-0.026)$ \\
$16$--$18$ (transition) & block $18$ & $0.712$ & $0.708\;(+0.004)$ & $0.815\;(-0.103)$ \\
$19$--$26$ (late)       & block $27$ & $0.537$ & $0.534\;(+0.003)$ & $0.539\;(-0.002)$ \\
\bottomrule
\end{tabular}
\end{table}

Three readings, in order of confidence:
\begin{itemize}
\item \textbf{Middle band: attention is causally necessary.} Removing
attention collapses truth from $0.793$ to $0.487$ (drop $+0.306$); removing
the FFN does nothing ($-0.026$). Attention does not merely read the signal
here, it builds it.
\item \textbf{Transition band: the FFN \emph{degrades} the signal.} Removing
attention again does nothing ($+0.004$), but removing the FFN \emph{raises}
truth AUC from $0.712$ to $0.815$ (drop $-0.103$). From layer $16$ onward the
FFN injects content that lowers true/false separability. This is the first
causal support for an earlier conjecture, that the FFN adds context which
amplifies the overlap on uncertain facts, though it rests on one band and one
dataset and should be treated as a lead, not a result.
\item \textbf{Late band: nothing remains to destroy.} The intact AUC is
already about $0.54$ (near chance): the truth signal has evaporated toward
the output, so neither ablation has an effect.
\end{itemize}

\textbf{This falsifies the initial ``attention builds, FFN consolidates''
hypothesis.} There is no late consolidation; the FFN's apparent late
\emph{readability} dominance (Sec.~\ref{sec:alternation}) was the readability
of a decaying signal, not active maintenance. The corrected picture: attention
constructs the truth representation in the middle layers; the FFN begins to
erode it in the transition band; by the late layers it is gone, as the model
shifts from representing the fact to predicting the next token.

\section{Discussion: what a single axis can and cannot read}
\label{sec:discussion}

It is tempting to view a cheap geometric probe as a low-cost version of
dictionary learning, reading superposition without expanding it. We argue
this is impossible, and that the impossibility is a capacity statement, not
an implementation gap. A projection is an inner product: it collapses a
$d$-dimensional state to a single scalar, discarding $d - 1$ dimensions. From
that scalar one cannot reconstruct which features were present, only how much
a chosen direction was. ``Reading the superposition'' requires asking the
question once per dictionary atom, which over an overcomplete dictionary of
size $m$ is exactly the $O(d\,m)$ sparse autoencoder. Compression (project to
a scalar) and disambiguation (expand to a sparse code) are not inverse
cheap-vs-expensive versions of one operation; they are different operations,
and only the second reads superposition. The gaps of
Sec.~\ref{sec:dim}, $0.062$ on homogeneous unknown material and $0.203$ on
the heterogeneous mix, are direct measurements of the truth information that
a single axis structurally cannot return.

The consequence is a clear scope. This method's value is \emph{structural and
cheap}: a training-free axis that, on facts the model knows, rivals a trained
probe at two dot products per token; a quantitative law relating the
readability of truth to the model's knowledge of the fact; and a route to the
polarity direction that needs no polarity labels. It is not a substitute for
sparse autoencoders and should not be presented as one.

\section{The axis reads truth and confidence}
\label{sec:meta}

The seven experiments establish that there is no second coordinate of truth
to find. This has a constructive consequence that is easy to miss, so we
state it plainly. The same axis admits \emph{two} readings, not one.

Its \emph{sign} classifies a statement: $\mathrm{Re} > 0$ is the true side,
$\mathrm{Re} < 0$ the false side. Its \emph{separability}, how cleanly true
and false split on a given set of facts, measures how well the model knows
that material. The geometry that answers ``is this statement true?''
therefore also answers, indirectly, ``does the model know this?'', because a
model that does not know produces overlapping rather than separated
representations (Figs.~\ref{fig:plane},~\ref{fig:plane_tqa}). These are not
two axes; they are two readings of one axis: the direction of the arrow, and
the sharpness of the arrow.

The reading is quantitative at the level of a distribution. The position-axis
separation is $0.795$ on CounterFact and $0.656$ on TruthfulQA: by this
measure the model ``knows'' the CounterFact material better than the
TruthfulQA material, read off the geometry alone, with no external grading.
We report this as a property of the dataset, not yet of the single fact.
Calibrating a per-fact confidence, a local separability around one statement
that predicts whether the model actually errs on it, is the natural next
step, and requires a behavioral label (does the model answer the open
question correctly?) that the minimal-pair design does not contain. We mark
it as future work rather than a claim. Section~\ref{sec:future} returns to
this label and to a prerequisite it exposed.

\section{Why unknown facts overlap: the signature of superposition}
\label{sec:why}

It remains to interpret \emph{why} the classes overlap on unknown facts
rather than simply asking \emph{that} they do. The overlap is not measurement
noise; it is the observable signature of unresolved superposition, and the
explanation follows from the capacity argument of Sec.~\ref{sec:discussion}.

A projection onto a single direction returns a scalar: it collapses the
$d$-dimensional state to one number and discards $d - 1$ dimensions. When the
model knows a fact, its truth content has been allocated a clean, dominant
direction of its own (the residual stream has, in effect, dedicated an axis
to it), and the single projection recovers almost all of it (position AUC
$0.795$, near the full-probe ceiling). When the model does not know a fact,
there is no dedicated direction: the weak truth-relevant signal remains
distributed across many superposed features, and no single axis can pull it
out of the tangle. Recovering it would require asking the question once per
dictionary atom, the expansion that a sparse autoencoder performs at cost
$O(d\,m)$, which is a different operation from a single projection, not a
cheaper version of it.

The consequence is exactly what the visual observations show. On known facts
the truth occupies its own direction, the classes form two regions, and the
centroid ratio is high. On unknown facts the truth is spread across the
superposition, the classes overlap, and the centroid ratio falls below one.
The dense, uniform overlap of true and false on unknown facts is therefore
not a failure of the probe but a direct, visual readout of superposition that
the probe structurally cannot resolve. The probe gaps of
Sec.~\ref{sec:dim} and the low centroid ratio measured here are views of the
same quantity: the truth information that lives off-axis on facts the model
does not know ($0.062$), inflated further when heterogeneous material forces
one axis to serve many truth frames at once ($0.203$). Subsequent work
Part~II begins to resolve that superposed
remainder into named, signed components: what this paper can only call
unresolved superposition acquires, there, a geometry.

\section{A note for future work: verify behavior before geometry}
\label{sec:future}

Two natural extensions, a per-fact behavioral confidence label
(Sec.~\ref{sec:meta}) and a study of the geometry of \emph{logical} (as
opposed to semantic) truth, share a prerequisite that is easy to overlook,
and we record it here as methodological guidance rather than as a result.

\paragraph{An informal behavioral check (not a controlled experiment).}
Before building any geometric study of a capability, one should confirm that
the capability exists behaviorally. We ran a small, qualitative probe of the
base model in open-ended completion (greedy decoding, no context carried
between prompts). Two things were apparent on a handful of prompts. First,
the model reproduces the common semantic facts used throughout this paper
(capitals, chemical symbols, authorship, currencies), consistent with the
``facts the model knows'' assumption behind the curated pairs. Second, on
symbolic transitive syllogisms built from nonce words (``a \emph{grok} is
heavier than a \emph{blarn}; a \emph{blarn} is heavier than a \emph{fendu};
therefore $\ldots$''), the base $1.5$B model does \emph{not} perform the
deduction: it deflects into multiple-choice or code scaffolding rather than
completing the transitive conclusion. We stress that this is an informal
observation on a few prompts, not a measured result, and we draw no
quantitative claim from it.

\paragraph{The principle.} Studying the geometry of a capability presupposes
the capability. Measuring the hidden-state geometry of a task the model
cannot perform yields noise, and that noise is easily mistaken for ``the
property is distributed/complex'' when the truthful reading is ``the property
is absent.'' A cheap behavioral feasibility check should therefore precede
any geometric study of \emph{logical} truth; on this model that study is not
feasible, which is itself useful to know in advance. The same principle
applies to transfer: before studying the geometry of truth in a second model,
one should first verify behaviorally that the second model knows the facts in
question. (In the subsequent multi-model work of
Part~II this principle is applied systematically: the
behavioral know-rate is measured before any geometry.)

\paragraph{A hypothesis to test, not a claim.} The way the base model
deflects on uncertain input, falling back on memorized \emph{format}
scaffolding (exam layouts, question/answer templates, code), suggests, but
does not establish, that part of the overlap on unknown facts may carry a
sentence-format component rather than pure superposition noise. This is
consistent with the format confound that makes mixing datasets inadvisable
(Sec.~\ref{sec:limits}). We flag it as a direction for a future, properly
controlled experiment, not as evidence.


\section{Limitations}
\label{sec:limits}

These results are confined to a single, small model (Qwen2.5-1.5B); the
entire battery, the positive results and all seven falsification experiments,
should be re-run at larger scale before the one-dimensionality is treated as
general (the partial replications available so far are summarized, with their
scope, in Part~II). It is plausible that departments,
dynamics, or a genuine second dimension emerge with scale, precisely because
a $1.5$B residual stream may be too compact to separate them. The datasets
are small (tens to a few hundred pairs), so confidence intervals are wide
despite the permutation control. Cross-validation is over pairs, not over
topics: pairs from the same CounterFact relation may fall in different folds,
so generalization to entirely unseen topics is not certified by the protocol
of Sec.~\ref{sec:protocol}; the within- vs.\ cross-category transfer
measurements of subsequent work Part~II
address the question directly. The ``unknown-fact'' interpretation of the
CounterFact gap and of the low separability is plausible but confounded: a
fact the model does not know and a fact whose truth is genuinely
multidimensional are not separated here, and the low-separability reading is
further confounded with sentence-format effects, which is why mixing datasets
must be avoided. Disentangling ignorance from genuine multidimensionality
(for example by filtering CounterFact to facts the model can answer, using
the behavioral label of Sec.~\ref{sec:future}) is necessary to make the
dimensionality claim quantitative. The anatomy results are Tier 2: margins in
the expanded-space and transition-band analyses are modest and based on
single datasets, the factual-knowledge label is noisy, and ablation measures
necessity for \emph{linear readability} of truth, not for the model's
behavior. We use last-token pooling; other aggregations may shift the axis.

\section{Summary of Part I}

Stripped of a decorative complex-plane framing, the positive results survive
on a real model under a falsifiable protocol: a training-free truth axis,
identified without labels up to one global sign, carries a permutation-tested
signal at cost $O(d)$; the gap between that axis and a full probe scales with
how well the model knows the fact, suggesting a knowledge-dependent
dimensionality of truth; and the difference SVD recovers the supervised
polarity direction of~\cite{burger2024} at cosine $0.959$ without polarity
labels. The one-dimensional reading withstands seven falsification
experiments, spanning a second dimension, depth-global structure, dynamics,
departments, energy, phase, and three-dimensional clustering, none of which
exhibits a richer separating coordinate. Beneath the axis, an exact
decomposition of the residual stream locates its origin: attention builds the
signal in the middle layers, causally, while the FFN erodes it after the
peak, and by the late layers the model has moved on from representing the
fact to predicting the next token. We read a result that resists seven
falsifications as strengthened, not weakened, interpret the residual overlap
on unknown facts as the visual signature of unresolved superposition, and
regard the axis as a light, honest, almost-free structural probe to sit
alongside generation. Natural next steps, partly realized in the subsequent
work summarized in Part~II, are the multi-model study
of the dimensionality law, a knowledge-controlled dataset built from a
behavioral label, and the category-level structure of the axis, each preceded
by a behavioral feasibility check of the kind described in
Sec.~\ref{sec:future}. All experiments in this manuscript are fully reproducible; the exact execution commands for both parts are consolidated at the end of the text.


\part*{Part II: Mechanics and Geometry Across Two Families}
\section{Scope and evidence levels}
\label{b:sec:scope}

Results are organized in three evidential levels and never mixed.
\textbf{Level A}: replicated across models with predictions written before
the runs. \textbf{Level B}: causal but on fewer models, or with multiple
independent methods on one. \textbf{Level C}: leads, declared as such
(Sec.~\ref{b:sec:leads}). Retractions and failed predictions are reported in
the text at the point where they occur; the largest, the correction of our
own flip interpretation, has its own section (\ref{b:sec:correction}).

\section{Setup}
\label{b:sec:setup}

Four base (non-instruct) models across two families:

\begin{table}[H]
\centering
\caption{Models. The truth peak $p$ is the block whose intact residual
maximizes held-out truth AUC (protocol of (Part~I)); ``block $b$''
denotes the output of transformer block $b$. Know-rate: fraction of $200$
CounterFact prompts whose true target the model reproduces under greedy
decoding (seed $0$), measured before any geometry. Axis and probe AUC:
CounterFact, $250$ pairs, held-out; gap $=$ probe $-$ axis.}
\label{b:tab:models}
\begin{tabular}{lccccccc}
\toprule
model & blocks & $d_{\mathrm{model}}$ & peak $p$ & know-rate & axis & probe & gap \\
\midrule
Qwen2.5-1.5B  & $28$ & $1536$ & $15$ & $30.0\%$ ($61/200$) & $0.792$ & $0.854$ & $0.062$ \\
Qwen2.5-3B    & $36$ & $2048$ & $16$ & $37.0\%$ ($74/200$) & $0.832$ & $0.881$ & $0.049$ \\
Llama-3.2-1B  & $16$ & $2048$ & $7$  & $35.5\%$ ($71/200$) & $0.690$ & $0.831$ & $0.141$ \\
Llama-3.2-3B  & $28$ & $3072$ & $9$  & $41.0\%$ ($82/200$) & $0.785$ & $0.897$ & $0.112$ \\
\bottomrule
\end{tabular}
\end{table}

Two regularities and one boundary, extending the dimensionality claim
of (Part~I): within each family the gap decreases as the know-rate
increases; across families a family-level offset remains, with Qwen
concentrating more of the truth signal on the single axis at comparable
knowledge. An earlier version of this table imported the Qwen2.5-1.5B row
from the first part's mixed-dataset column while the other rows were
measured on CounterFact alone, a protocol inhomogeneity caught by
re-measurement; all rows above come from one command under one protocol,
and the earlier phrasing that the gap ``tracks knowledge every time'' is
retired in favor of the two-part statement.

Dataset: \texttt{NeelNanda/counterfact-tracing} (pinned revision
\texttt{c945b08\ldots}, loaded as Parquet only), $250$ minimal pairs unless
stated. Decomposition runs use \texttt{float32} and are gated by the
additive identity check of (Part~I) (median relative errors
$0.65$--$1.23 \times 10^{-7}$ across the runs of this paper);
extraction-only runs use \texttt{bfloat16}. The spread of the four peaks in
relative depth ($0.32$--$0.54$) already anticipates the negative of
Sec.~\ref{b:sec:negatives}: no fixed depth, absolute or relative, locates
truth.

\section{Methods}
\label{b:sec:methods}

\subsection{Frames and the fixed-frame contribution statistic}
\label{b:sec:fixedaxis}

Within a block the residual stream updates additively,
$h_{b+1} = h_b + a_b + f_b$, with $a_b$ the attention contribution and $f_b$
the FFN contribution, captured by forward hooks (Part~I).

\begin{definition}[Truth frame]
The \emph{post frame} $F_b^{\mathrm{post}}$ is the oriented truth axis $v_1$
fit (on training folds only, true side positive) on the block output
$h_{b+1} = h_b + a_b + f_b$; it contains $f_b$ and $a_b$. The \emph{pre
frame} $F_b^{\mathrm{pre}}$ is the same construction on the pre-FFN state
$h_b + a_b$; it contains $a_b$ but excludes $f_b$. A frame \emph{exists} at
$b$ when the axis fit there clears its permutation null held-out.
\end{definition}

The statistic for a contribution $c$ at layer $L$ measured against a frame
$F$ with axis $v_1$ is the held-out class gap and its effect size,
\begin{equation}
\mathrm{gap}(c; F) \;=\; \overline{(v_1 \cdot c)}\big|_{\mathrm{true}}
- \overline{(v_1 \cdot c)}\big|_{\mathrm{false}},
\qquad
d'(c; F) \;=\; \frac{\mathrm{gap}(c; F)}{\text{pooled std}},
\end{equation}
with the pooled standard deviation over the held-out projections of both
classes. A negative FFN gap means the FFN literally writes toward the false
side of the frame. Stability criterion, fixed in code before the
consolidation runs: a value is \emph{stable} if its magnitude exceeds both
$2\sigma$ across $5$ cross-validation seeds and a floor of $0.03$.

\subsection{The exact SwiGLU split}
\label{b:sec:swiglu}

The FFN computes $f = W_{\mathrm{down}}(g \odot u)$ with
$g = \mathrm{silu}(W_{\mathrm{gate}} x)$, $u = W_{\mathrm{up}} x$. The
projection of $f$ on $v_1$ is $w \cdot (g \odot u)$ with
$w = W_{\mathrm{down}}^{\top} v_1$, the residual axis pulled back through
the down-projection; no axis is ever fit in the expanded space, avoiding the
$p \gg n$ regime (more dimensions than samples) entirely. With the symmetric
intra-pair split ($\Delta g = g_t - g_f$, $\bar g = (g_t + g_f)/2$, likewise
for $u$) the mixed term cancels and the identity is exact:
\begin{equation}
w \cdot (g_t \odot u_t - g_f \odot u_f)
= \underbrace{w \cdot (\Delta g \odot \bar u)}_{\text{gate term}}
+ \underbrace{w \cdot (\bar g \odot \Delta u)}_{\text{value term}}.
\label{b:eq:split}
\end{equation}
Verified numerically at $7.4 \times 10^{-8}$ relative error on real data.

\subsection{Interventions}
\label{b:sec:interventions}

Ablating a component means replacing its additive contribution ($f_{\ell}$
or $a_{\ell}$) with the zero vector through a forward hook, for every layer
$\ell$ in a band (peak$+1$ to peak$+3$; readout at the band end, on every
model), during the forward pass: the weights are untouched, the module still
computes and its output is discarded, and the residual update reduces to
$h_{\ell+1} = h_{\ell} + a_{\ell}$ (for FFN ablation) or
$h_{\ell+1} = h_{\ell} + f_{\ell}$ (for attention ablation). Two
refinements. (i) A
\emph{frozen control} obtained for free: with both components zeroed over
the band, the readout state is identically the band-entry state, so the
frozen reference is read off the intact pass. (ii) \emph{Gate-freeze and
value-freeze}: at the reading position, $g$ (respectively $u$) is replaced
by its intra-pair mean cached from the intact run, killing the class-driven
variation of one stream while the FFN stays on.

\paragraph{On the validity of the freezes.} Any activation-level
intervention risks pushing the network off its data manifold, so three
properties of this design matter. The replacement value is the
\emph{intra-pair mean}: a convex combination of the two real activations
of two nearly identical inputs, which is as close to on-manifold as an
intervention of this kind gets. The two freezes are symmetric by
construction, so their \emph{contrast} (value recovery $+0.064$ to
$+0.123$ against gate recovery $+0.008$ to $+0.056$ across the four
models) controls for generic disruption: an off-manifold artifact has no
reason to respect the gate/value distinction. And the interventional
verdict agrees with the purely observational exact split of
Eq.~\ref{b:eq:split}, which touches nothing. The residual risk is
declared, not dismissed.


\subsection{Static circuit norms}
\label{b:sec:circuits}

Per attention head (GQA-aware), the gauge-invariant circuit norms
$\lVert W_q^{\top} W_k \rVert_F$ (routing) and $\lVert W_o W_v \rVert_F$
(writing)~\cite{elhage2021}. Individual weight norms are gauge-dependent
($W_q \to \alpha W_q$, $W_k \to W_k / \alpha$ leaves the function
unchanged); the products are not. A gauge self-test gates the computation.
The \emph{crossover} is the first layer from which the median-normalized
OV/QK ratio stays above $1$ for at least three consecutive layers, a
criterion fixed in code before any run.

\subsection{The axis-provenance control}
\label{b:sec:provenance}

The objection this method answers, raised by the author against the
author's own instrument, is that a result anchored to ``the block where the
ruler was fit, plus one'' could be a property of the ruler rather than of
the FFN. The control: fit frames at several blocks
$b \in \{p{-}4, p{-}2, p, p{+}2\}$, both post and pre, and recompute the
full $d'$ profile of each component under each frame. The discrimination
criterion was fixed before the runs and validated on synthetic ground
truth: a planted global axis with a planted writer flip stays pinned under
every frame; a planted circular writer collapses under the pre frame
(collapse ratio $0.04$) while planted genuine alignment survives ($0.99$).
The released tools print full matrices with no automatic verdict; the
reason is a repeated methodological lesson recorded in
Sec.~\ref{b:sec:negatives}.

\subsection{Per-category axes, transfer, decoding, and the Mantel statistic}
\label{b:sec:catmethods}

CounterFact is grouped by Wikidata relation; the top-$K$ relations by
unique pairs are used ($K{=}8$ for the registered falsification, $60$ pairs
each): P103 mother tongue, P1412 languages spoken, P176 manufacturer, P27
citizenship, P30 continent, P37 official language, P413 sport position,
P495 country of origin. Per-category truth axes are fit at the peak block,
each oriented within its own category, so the sign of a cross-category
cosine is information: negative means the shared direction reads the other
category's truth backwards. Four measurements: (i) signed cosine matrices
at the peak and at an early block (block $2$, the lexical surface control);
(ii) the held-out \emph{transfer matrix}, cell $(i,j)$ being the held-out
AUC of category $i$'s axis on category $j$'s pairs; (iii) nearest-centroid
\emph{decoding} of the category from the FFN's class-signed write direction
$\Delta f$ at peak$+1$, against the same decoder on early residual
differences (the lexical baseline); (iv) the \emph{Mantel test}, a
permutation test for the correlation between two distance matrices, here
applied under category-label permutation to compare the two families'
signed-cosine arrangements.

\section{First movement: the flip, its correction, and the relational law}
\label{b:sec:flip}

\subsection{The peak-relative flip (the measurement that survives)}

For each model the prediction was written before the run: under the
peak-fit post frame, the FFN's gap is pro-truth at the peak block $p$ and
flips to stable anti-truth at $p{+}1$.

\begin{table}[H]
\centering
\caption{The flip at peak$+1$ under the peak-fit post frame. CounterFact
$250$ pairs, $5$ seeds, permutation null with $100$ within-pair swaps
($\hat p = 0.0099$ on every row). All four rows from the uniform July 10
consolidation; the Qwen2.5-1.5B row supersedes the earlier single-seed scan
($+0.728$/$-0.692$), with which it agrees to the second decimal.}
\label{b:tab:flip}
\begin{tabular}{lcccc}
\toprule
model & $\mathrm{gap}_{\mathrm{ffn}}(p)$ & $\mathrm{gap}_{\mathrm{ffn}}(p{+}1)$ & $d'(p)$ & $d'(p{+}1)$ \\
\midrule
Qwen2.5-1.5B & $+0.730 \pm 0.002$ & $-0.691 \pm 0.002$ & $+0.58$ & $-0.82$ \\
Qwen2.5-3B   & $+0.575 \pm 0.005$ & $-0.926 \pm 0.002$ & $+0.68$ & $-1.25$ \\
Llama-3.2-1B & $+0.092 \pm 0.001$ & $-0.111 \pm 0.001$ & $+0.54$ & $-0.94$ \\
Llama-3.2-3B & $+0.087 \pm 0.001$ & $-0.167 \pm 0.001$ & $+0.49$ & $-1.26$ \\
\bottomrule
\end{tabular}
\end{table}

The four peaks sit at different absolute and relative depths
(Table~\ref{b:tab:models}); the flip sits at peak$+1$ in all four. Aligned in
peak-relative coordinates and normalized by the flip amplitude, the
per-layer profiles collapse onto a common shape: rise into the peak, spike
at $+1$, decay toward zero (Fig.~\ref{b:fig:collapse}). The measurement is
correct and reproduces; its interpretation, that the flip is anchored to
the peak, is what the next subsection tests.

\subsection{The correction: the control bites its owner}
\label{b:sec:correction}

Under frames fit at $p{-}4$, $p{-}2$, $p$, $p{+}2$
(Sec.~\ref{b:sec:provenance}), the transition does not stay at $p{+}1$: it
tracks the frame. On every model the post-frame diagonal reads pro at $b$
and anti at $b{+}1$ for every $b$ where a frame exists
(Table~\ref{b:tab:post}). The peak-anchored flip was the $b{=}p$ slice of a
relational pattern present at every such block. An accidental replication
is worth recording: a consolidation run launched with a wrong default
re-measured Qwen2.5-3B with the frame at block $15$, off its peak, and
returned the same pro/anti pair at $15/16$ ($+0.541$/$-0.545$, $5$ seeds,
$\hat p = 0.0099$); the accident confirms the law at a frame no one chose.

\begin{table}[H]
\centering
\caption{Post-frame diagonals: per fit block $b$, the FFN's held-out $d'$
at $b$ and at $b{+}1$. Pro at $b$, anti at $b{+}1$, wherever a frame
exists.}
\label{b:tab:post}
\begin{tabular}{lccc}
\toprule
model (peak) & frame $b$ & $d'_{\mathrm{ffn}}(b)$ & $d'_{\mathrm{ffn}}(b{+}1)$ \\
\midrule
Qwen2.5-1.5B ($15$) & $13$ / $15$ / $17$ & $+0.61$ / $+0.57$ / $+0.52$ & $-0.79$ / $-0.82$ / $-0.73$ \\
Qwen2.5-3B ($16$)   & $14$ / $16$ / $18$ & $+0.06$ / $+0.69$ / $+0.73$ & $-0.09$ / $-1.23$ / $-0.79$ \\
Llama-3.2-1B ($7$)  & $5$ / $7$ / $9$    & $+0.13$ / $+0.54$ / $+0.19$ & $-0.07$ / $-0.95$ / $-0.54$ \\
Llama-3.2-3B ($9$)  & $7$ / $9$ / $11$   & $+0.37$ / $+0.49$ / $+0.57$ & $-0.42$ / $-1.25$ / $-1.39$ \\
\bottomrule
\end{tabular}
\end{table}

Under the pre frame the FFN's ``pro at own block'' term does not merely
shrink; it inverts sign (Table~\ref{b:tab:pre}). Measured against the frame
that exists at the moment it writes, the FFN is strongly anti-aligned. The
pro term of the post frame was, in substantial part, the frame measuring
its own ingredient: $f_b$ is inside $h_{b+1}$ and drags the axis fit there
toward itself. Where the axis is weak or absent, the pattern correctly
vanishes: frames fit before the Qwen-3B ignition (fit-block cosines to the
peak axis $0.049$ and $0.079$) and before the Llama-1B peak return noise.
What the failure bought: a general law instead of a special case, and a
live demonstration that the safeguards bind their owner.

\subsection{The relational law}
\label{b:sec:law}

\begin{table}[H]
\centering
\caption{Pre-frame diagonals: the same cells with $f_b$ excluded from the
frame's data. The FFN inverts sign; attention, measured identically, is the
mirror image.}
\label{b:tab:pre}
\begin{tabular}{lccc}
\toprule
model & frame $b$ & $d'_{\mathrm{ffn}}(b)$ & $d'_{\mathrm{attn}}(b)$ \\
\midrule
Qwen2.5-1.5B & $13$ / $15$ / $17$ & $-0.63$ / $-0.41$ / $-0.78$ & $+1.56$ / $+1.31$ / $+1.04$ \\
Qwen2.5-3B   & $14$ / $16$ / $18$ & $-0.14$ / $-1.16$ / $-0.46$ & $+0.26$ / $+1.38$ / $+1.40$ \\
Llama-3.2-1B & $7$ / $9$          & $-1.17$ / $-0.82$           & $+1.26$ / $+0.82$ \\
Llama-3.2-3B & $7$ / $9$ / $11$   & $-0.81$ / $-1.43$ / $-1.41$ & $+0.99$ / $+1.63$ / $+1.75$ \\
\bottomrule
\end{tabular}
\end{table}

The decisive zone is the sub-diagonal of the pre matrices: contributions at
layer $L$ measured against frames of blocks strictly before $L$, which
contain neither $f_L$ nor $a_L$ and are clean of every containment effect.
There (Fig.~\ref{b:fig:asymmetry}) the FFN is negative and attention positive
nearly everywhere a frame exists, on all four models. At the peak, under
the same clean frame, the two components are near mirror images:
$+1.31/-0.41$ (Qwen-1.5B), $+1.38/-1.16$ (Qwen-3B), $+1.26/-1.17$
(Llama-1B), $+1.63/-1.43$ (Llama-3B). Pre-registered predictions for the
two Llamas, written before those runs, were all confirmed. The regularity
can be stated compactly.

\begin{law}[Propagate vs.\ overwrite]
\label{b:law:relational}
Let $F$ be a truth frame that exists at block $b$, and let $L \geq b$ be a
layer whose contributions are not contained in $F$ (any $L$ for pre frames
at $b = L$; any $L > b$ for both frame types). Then, as an empirical
regularity across the four models studied: $d'(f_L; F) < 0$ and
$d'(a_L; F) > 0$ wherever both are stable. Attention propagates frames it
did not write; the FFN opposes them and deposits the material of later
frames.
\end{law}

The law dissolves the earlier implicit picture of an axis being eroded like
an object: the axis is not an object in the model but a moving statistical
summary of where the class difference lives, and the FFN's class-signed
writes are what successive frames are made of. This is the dynamical face
of the key-value write-head physiology attributed to
MLPs~\cite{meng2022}, and a candidate microscopic writer for the
macroscopic finding of~\cite{rotational2026}.

\begin{figure}[t]
\centering
% Upload flip_collapse.png alongside the .tex.
\includegraphics[width=0.9\textwidth]{flip_collapse.png}
\caption{The FFN's frame writing ($d'$) across depth, aligned at the truth
peak, four models. Right: normalized by the flip amplitude.}
\label{b:fig:collapse}
\end{figure}

\begin{figure}[t]
\centering
% Upload provenance_asymmetry.png alongside the .tex.
\includegraphics[width=0.95\textwidth]{provenance_asymmetry.png}
\caption{Propagate vs.\ overwrite on four models: held-out $d'$ of each
component's contribution (FFN above, attention below) against pre frames
fit at four blocks. The FFN contribution is negative wherever a frame exists; the
attention contribution is positive.}
\label{b:fig:asymmetry}
\end{figure}

\section{The erosion: real, frame-free, and carried by the value stream on four models}
\label{b:sec:erosion}

The correction separates two phenomena the peak-anchored reading conflated:
frame-relative rotation (what the flip was) and frame-free decay of class
signal after the peak (the genuine erosion). The erosion stands on
measurements that use no frozen frame, now on all four models
(Table~\ref{b:tab:ablate}); the predictions for the Llama rows were
registered before those runs and confirmed. This section is Level A.

\begin{table}[H]
\centering
\caption{Causal ablation over the band peak$+1$ to peak$+3$, readout at
band end; held-out refit truth AUC. FFN-off restores signal on every model;
attention-off is neutral on the smaller models and destructive on both 3B
models (the tug-of-war of Sec.~\ref{b:sec:traits}).}
\label{b:tab:ablate}
\begin{tabular}{lcccc}
\toprule
model & intact & FFN-off & attn-off & frozen \\
\midrule
Qwen2.5-1.5B & $0.712$ & $0.815$ & $0.708$ & $0.793$ \\
Qwen2.5-3B   & $0.793$ & $0.846$ & $0.738$ & $0.831$ \\
Llama-3.2-1B & $0.551$ & $0.744$ & $0.531$ & $0.691$ \\
Llama-3.2-3B & $0.740$ & $0.847$ & $0.593$ & $0.785$ \\
\bottomrule
\end{tabular}
\end{table}

\begin{table}[H]
\centering
\caption{Inside the SwiGLU, four models. Left: the exact split of
Eq.~\ref{b:eq:split} at the flip layer (peak-fit frame). Right: causal
freezes over the band (AUC recovery vs.\ intact). The anti-truth write and
its causal carrier are the value stream everywhere; the gate is neutral or
mildly pro-truth.}
\label{b:tab:swiglu}
\begin{tabular}{lccccc}
\toprule
 & \multicolumn{3}{c}{split at flip} & \multicolumn{2}{c}{freeze recovery} \\
model & gate & value & total & gate & value \\
\midrule
Qwen2.5-1.5B & $+0.041$ & $-0.734$ & $-0.692$ & $+0.040$ & $+0.089$ \\
Qwen2.5-3B   & $-0.120$ & $-0.806$ & $-0.926$ & $+0.008$ & $+0.064$ \\
Llama-3.2-1B & $+0.016$ & $-0.127$ & $-0.110$ & $+0.028$ & $+0.123$ \\
Llama-3.2-3B & $+0.021$ & $-0.186$ & $-0.166$ & $+0.056$ & $+0.122$ \\
\bottomrule
\end{tabular}
\end{table}

Three annotations keep the table honest. First, the earlier hypothesis that
the gate carries the erosion (which neurons open switching the stream
toward next-token context) is falsified in the same direction on every
model; at the peak, by contrast, the split shows the gate carrying about
two thirds of the pro-truth construction (Qwen-1.5B: gate $+0.486$, value
$+0.242$). Second, on Llama-3.2-3B the value-freeze lands at $0.862$, above
the frozen reference $0.785$: with the value stream's class variation
silenced, the band's attention adds signal, which is the tug-of-war of the
next section seen causally. Third, the earlier lead ``FFN-off beats
frozen'', demoted to noise on the Qwen models ($+0.022$/$+0.015$, both
under the stability floor), resurfaces above the floor on both Llamas
($+0.052$/$+0.062$); the demotion was correct for Qwen, and the effect is
recorded as family-dependent, without a dedicated null, hence an
observation and not a claim. A frame-free rotation check closes the
projection-artifact objection on three models: the cosine of the mean FFN
pair-difference with $v_1$ jumps sign in a single step at peak$+1$ with the
write magnitude intact (deltas $-0.741$, $-0.654$, $-0.632$ on Qwen-1.5B,
Qwen-3B, Llama-1B), where genuine rotation predicts a smooth multi-layer
sweep.

\section{Two traits, discriminated: tug-of-war is scale, ignition is family}
\label{b:sec:traits}

Two pre-registered discriminations with opposite outcomes, which together
show that scale and training recipe leave distinct, measurable signatures.

\paragraph{Tug-of-war is a trait of scale.} In the erosion band, removing
attention is neutral on the two smaller models (refit change $-0.004$ on
Qwen-1.5B, $-0.020$ on Llama-1B, both under the floor) and destructive on
both 3B models ($-0.055$ on Qwen-3B, $-0.147$ on Llama-3B). The registered
discrimination read: present on Llama-3B and absent on Llama-1B implies
scale; absent on both implies a Qwen-family trait. Outcome: scale. In the
larger models the post-peak band is a tug-of-war in which attention still
writes pro-truth while the FFN's value stream overwrites, not a passive
decline.

\paragraph{Ignition is a trait of the Qwen family.} The Qwen-3B per-layer
AUC curve switches on in a single layer ($0.526 \to 0.790$), driven by a
coordinated two-layer construction burst; the Qwen-1.5B curve is a smooth
ramp. The registered discrimination: scale predicts a step in Llama-3.2-3B;
family predicts a ramp. Measured: a ramp ($0.526 \to 0.657 \to 0.664 \to
0.749 \to 0.785$ across levels $6$--$10$). Ignition is, on present
evidence, a trait of the Qwen recipe at the 3B scale, not of scale itself.

\section{A pre-registered negative, cross-family: the static crossover does not track the peak}
\label{b:sec:negatives}

Motivated by an external observation that value-to-query weight norms scale
with depth~\cite{headgenome}, the gauge-invariant version was tested: does
the layer where the writing circuit overtakes the routing circuit predict
the truth peak? The prediction for the Llama models, registered before
their runs, was that it does not. Outcome (Table~\ref{b:tab:crossover}): the
crossover sits at a stable late relative depth on all four models
($0.61$--$0.75$) while the peak wanders ($0.32$--$0.54$); the two coincide
only on Qwen2.5-1.5B, the accident that spawned the hypothesis the first
working note falsified across scale. The depth-scaling of the circuits may
be real and general; it is not where truth lives. The late-layer OV
explosion (Qwen-1.5B sums $338.9$/$322.0$ at layers $26$--$27$) instead
matches the decoupling of Sec.~\ref{b:sec:leads}.

\begin{table}[H]
\centering
\caption{The OV/QK crossover (first layer sustaining ratio $>1$ for three
consecutive layers) against the truth landmarks, four models. Distance is
to the nearer of peak and flip.}
\label{b:tab:crossover}
\begin{tabular}{lcccc}
\toprule
model & crossover & peak / flip & distance & rel.\ depth (cross vs.\ peak) \\
\midrule
Qwen2.5-1.5B & $17$ & $15$ / $16$ & $1$ & $0.61$ vs.\ $0.54$ \\
Qwen2.5-3B   & $24$ & $16$ / $17$ & $7$ & $0.67$ vs.\ $0.44$ \\
Llama-3.2-1B & $12$ & $7$ / $8$   & $4$ & $0.75$ vs.\ $0.44$ \\
Llama-3.2-3B & $17$ & $9$ / $10$  & $7$ & $0.61$ vs.\ $0.32$ \\
\bottomrule
\end{tabular}
\end{table}

\paragraph{A methodological note on automatic verdicts.} During the
regeneration of this table, the circuit tool's hardcoded landmark
annotations (written for one model) mislabeled the truth peak on the other
three and printed a spurious coincidence verdict for Llama-3.2-3B. It was
the third automatic verdict to misread its own numbers in this project's
history, after two in the provenance tooling. All released tools now print
matrices and distances with no verdict prose; the reading belongs to the
researcher. The raw norm tables, which depend only on the weights, were
unaffected.

\section{Second movement: the axis is a semantically signed mixture}
\label{b:sec:mixture}

If the FFN writes truth along directions that depend on what kind of fact
is being written, the single axis is a mixture, and its ingredients can be
measured. Design and statistics as in Sec.~\ref{b:sec:catmethods}; the
registered falsification runs on the two 3B models. Llama-3.2-1B is
excluded from this section by a scoping rule stated in advance: it is the
weakest-axis model of the four (axis AUC $0.690$), and per-category axes
fit on $60$ pairs there risk reporting noise as geometry.

\paragraph{The signed geometry.} At the peak, the language-relation block
coheres (Qwen-3B cosines $+0.27$ to $+0.76$; Llama-3B $+0.55$ to $+0.69$),
continent anti-aligns with that block (Qwen $-0.07$ to $-0.20$; Llama
$-0.09$ to $-0.15$), sport-position is near-orthogonal to everything, and
country-of-origin sides with the block (strongly on Llama, $+0.51$ to
$+0.75$). The anti-alignment is confirmed by an independent statistic:
transfer from the continent axis into the language block falls below
chance (Llama: AUC $0.198$--$0.33$); on Qwen the below-chance transfer is
bidirectional, on Llama one-directional, a partial replication reported as
such. Within-category transfer exceeds cross-category transfer on both
models (Qwen $0.662$ vs.\ $0.596$; Llama $0.898$ vs.\ $0.808$), which also
bears on the topic-generalization question raised for the protocol of Part~I.

\paragraph{The reorganization with depth.} The early-block matrices are
not flat: they carry shared corpus statistics (the two families' early
matrices correlate at $r = +0.934$). Depth does not create arrangement from
nothing; it re-files it. Citizenship sides with geography at the embedding
level and with the language block at the peak, on both families. The
extreme case: the axes of languages-spoken and official-language start near
anti-parallel ($\cos = -0.90$ Qwen, $-0.86$ Llama, presumably a lexical
opposition in how the two templates use language tokens) and merge at the
peak ($+0.76$, $+0.62$). Within each model, the early and peak
arrangements are uncorrelated or anti-correlated ($r = -0.547$ Qwen,
$-0.175$ Llama): the peak structure contradicts, rather than inherits, the
lexical surface.

\paragraph{The FFN writes truth along category-dependent directions.}
Decoding the relation from the unit write direction $\Delta f$ at
peak$+1$: $74.17\%$ (Qwen-3B) and $81.67\%$ (Llama-3B) against chance
$12.5\%$ and permutation null $95$th percentiles near $15.5\%$
($\hat p = 0.0099$ both); lexical baseline $39.17\%$ and $42.08\%$. Read
together with Law~\ref{b:law:relational}: the component that overwrites truth
frames does so along category-dependent directions. The mixture is not
only in the representation but in the writing.

\paragraph{The registered falsification and its verdict.} Five predictions
were written before any Llama measurement of this geometry; outcomes in
full, partials included. \textbf{R1} (language block on Llama, cosines in
$[+0.15, +0.45]$): structure confirmed, magnitude outside the registered
window on the high side. \textbf{R2} (continent anti-aligned, below-chance
transfer both directions): signs confirmed; bidirectionality held only on
Qwen. \textbf{R3} ($\Delta f$ decoding at least $10$ points over lexical):
confirmed with a margin near $40$ points. \textbf{R4} (the citizenship
early-to-peak switch replicates): confirmed, plus the stronger merge.
\textbf{R5} (the arrangement law): confirmed; off-diagonal peak matrices,
Qwen-3B vs.\ Llama-3B, Pearson $r = +0.770$ (Spearman $+0.893$), Mantel
null $95$th percentile $+0.498$, $p = 0.0009$; robust to excluding the
anomalous relation P176 ($r = +0.808$, $p = 0.0037$).

\paragraph{The arrangement law.} Combined with the within-model
reorganization, the reading is: the two families share a lexical starting
arrangement, abandon it, and converge on the same functional arrangement.
Same start, same move, same destination, different machines. The
arrangement of the truth mixture is a property of the knowledge domain, not
of the training recipe: a signed, truth-specific instance of cross-model
representation convergence~\cite{platonic2024}. In hindsight this section
also explains a Part~I measurement from the other side: forcing
one axis onto heterogeneous material costs $0.203$ of held-out AUC against
a full probe, versus $0.062$ on homogeneous material (Part~I),
because the single axis is being asked to summarize a signed mixture with
one number.

\paragraph{Anomaly, declared.} On Llama-3B, relation P176 (manufacturer)
shows within-CV and incoming transfer of $1.000$ from every axis:
universal perfect separability, which suggests surface structure in that
relation's pairs rather than truth geometry. It is excluded in the R5
robustness check (which holds) and flagged for inspection before any claim
rests on it.

\begin{figure}[t]
\centering
% Upload categories_geometry.png alongside the .tex.
\includegraphics[width=0.95\textwidth]{categories_geometry.png}
\caption{The category geometry: early vs.\ peak signed-cosine matrices,
and decoding vs.\ the lexical baseline.}
\label{b:fig:catgeom}
\end{figure}

\begin{figure}[t]
\centering
% Upload arrangement_law.png alongside the .tex.
\includegraphics[width=0.95\textwidth]{arrangement_law.png}
\caption{The arrangement law (R5). Left: the two $K{=}8$ peak
signed-cosine matrices. Center: their off-diagonal cells against each
other. Right: cross-model agreement at the peak and at the embedding
level, and the within-model reorganization between them.}
\label{b:fig:law}
\end{figure}

\section{The measured dictionary and the sparse contract}
\label{b:sec:contract}

\texttt{dictionary\_export.py} saves, per model, the measured dictionary:
the per-category oriented axes, the global axis, and the FFN write
centroids, with signed matrices and full metadata in one loadable file (a
$[K \times d]$ bundle of tens of kilobytes): a $K$-direction truth and
category monitor whose baseline performance is quantified above. The
sparse autoencoder that this series hands to engineering is justified only
if, against matched nulls, it (1) aligns with the per-category axes, (2)
respects the semantic geography, (3) reproduces the signs, including the
inverted loading on continent and the depth reorganization, and (4) beats
the measured dictionary on class-signal reconstruction and category
decoding. Either outcome is content.

\section{Leads, declared as leads}
\label{b:sec:leads}

\paragraph{A second readable dimension in the mid-knowledge regime.} On
Llama-3.2-1B at the peak, a held-out logistic on $(\mathrm{Re},
\mathrm{Im})$ reaches $0.783$ vs.\ $0.690$ for the single axis; on
Llama-3.2-3B phase scores match or beat the axis in the decay zone. In the
Qwen models the second dimension never added more than $+0.02$. Read
through the dimensionality relation, this is the expected geometric
signature of the mid-knowledge regime: as knowledge decreases, off-axis
components become readable. It refines, and does not contradict, the
one-dimensionality claim of (Part~I), which was always scoped to
known facts. No dedicated null on the margin yet.

\paragraph{Per-head anatomy.} The attention contribution is a sum of heads
through the output projection, and the fixed-frame instrument applies
unchanged to each head's slice. Whether the pro-truth propagation at the
peak, and the 3B tug-of-war, are carried by a few identifiable heads or
distributed across the pack connects this anatomy
to~\cite{headgenome}.

\paragraph{Framing bias with an open lexical confound.} A $2{\times}2$
moral/prudential design on Qwen2.5-1.5B: framing does not rotate the axis
(transfer $0.979$--$0.99$), but there is a large level bias among
equally-true statements ($-2.75\sigma$ at layer $11$), and the SVD of
mixed-framing pairs recovers a framing direction at $|\cos| = 0.912$ in
PC2, the first part's recovery structure on a new factor. The
confound: the bias is already $-2.86\sigma$ at layer $1$, which points to
the lexical surface; crossed templates are required before this is more
than a lead.

\paragraph{Late-layer decoupling.} In all inspected models the total
volume written by the FFN grows toward the output while its projection on
the truth frame goes to zero and the intact truth AUC decays toward
chance. The terminal regime is a decoupling: the model writes enormously,
and none of it speaks of truth. Terminal sign changes in the gap are
zero-crossings of a dying signal, below both the stability floor and the
null, and must not be read as structure.

\paragraph{Frame-aware re-reading of the split.} The observational split
of Table~\ref{b:tab:swiglu} is measured under a peak-fit frame and should be
re-read against pre frames; the causal freezes are frame-free and the
attribution stands regardless.

\section{Limitations}
\label{b:sec:limits}

Four models, two families; a third family would promote the arrangement
law from law-candidate to law and widen the relational law's base. The pre
frame removes the direct containment circularity but not indirect paths
(accumulated earlier writes shape the pre frame). All measurements use
last-token states over one dataset (CounterFact, pinned revision) with one
statement form per relation; template syntax may contribute to the
continent anti-alignment, and crossed templates are the known cure. No
permutation nulls were run on the provenance matrices themselves; the
diagonal magnitudes sit far above the null scales measured on identical
statistics, but the formal check is future work. Ablation measures
necessity for linear readability of truth, not for behavior; the freezes
act at the reading position only. The category section uses $60$ pairs per
relation, no bootstrap error bars, and one degenerate relation on Llama
(declared). The know-rate is a string-match lower bound. The
propagate/overwrite reading is a description of measured alignments, not
yet a circuit-level mechanism; why attention inherits frames and the FFN
opposes them is the open question this paper sharpens rather than answers.




\section{Conclusion}

Two movements, one instrument, one discipline. The mechanics: attention
propagates truth frames it did not write; the FFN opposes the frame of its
moment and deposits the material of the next; the genuine post-peak erosion
survives every control and is carried, on all four models, by the SwiGLU
value stream. The geometry: the single axis is a semantically signed
mixture whose arrangement the model does not inherit from the lexical
surface but constructs, and constructs the same way across two families,
because the arrangement belongs to the knowledge, not to the machine.
Between the two movements sits the correction that makes them trustworthy:
a pre-registered claim of ours, confirmed on three models, was exposed as
the special case of a better law by a control built against our own
instrument, and the paper reports the failure with the same care as the
findings, because a result that has survived its author's own best attack is
worth more than one that has never been attacked. The handoff to the next
stage is concrete: a measured dictionary, a four-criterion contract for
any sparse autoencoder that claims to do better, and a per-head question
with its instrument already built.


\section*{Acknowledgments}

The author thanks Boaz Nadler for detailed and constructive comments on an
earlier version of Part~I, which prompted, among many improvements, the
formalization of its Method section and Lemma~\ref{lem:labels}; the
standards of definition and precision of those comments are applied
throughout this manuscript.

\section{Reproducibility, Part I}
\label{sec:repro}

Every number in this paper is reproduced by one command. \texttt{truth\_probe.py}
runs the controlled experiments (the positive results and falsification
experiments 1--4); \texttt{canvas.py} renders the visual observations
(experiments 5--7); \texttt{anatomy.py}, \texttt{anatomy\_consolidate.py},
\texttt{anatomy\_expanded.py}, and \texttt{ablation.py} produce the anatomy of
Sec.~\ref{sec:anatomy} and embed the identity, null, and stability safeguards
described there. Each result in \texttt{truth\_probe.py} is an \emph{isolated}
subcommand: running one task loads, extracts, and computes only its own data,
so a reviewer who wants to verify one figure runs exactly one command. All
scripts load the model with safetensors only and load external datasets as
Parquet only, never with remote code. A CUDA GPU is optional.

All scripts are archived on Zenodo (DOI 10.5281/zenodo.20938285) \\ and  mirrored
at \texttt{github.com/Francesco-Marhel/TruthProbe}.

\paragraph{Setup.}
\begin{verbatim}
pip install torch transformers datasets
\end{verbatim}

\paragraph{Positive results (Tables~\ref{tab:main},~\ref{tab:polarity}).}
\begin{verbatim}
# Table 1, top row (curated pairs): held-out AUC, ablation,
#                   full-probe baseline, permutation p-value
python truth_probe.py signal --dataset builtin --with-2d --baseline --perm 200

# Table 1, bottom row (CounterFact / TruthfulQA mix)
python truth_probe.py signal --dataset mix --max-pairs 250 \
       --with-2d --baseline --perm 200

# Table 2, top (affirmative -> negated flip; the 2D polarity structure)
python truth_probe.py polarity

# Table 2, bottom (SVD recovers the polarity direction t_P)
python truth_probe.py recovery
\end{verbatim}

\paragraph{Know-rate (Secs.~\ref{sec:data} and~Part~II) and the CounterFact row of Table~\ref{tab:main}.}
\begin{verbatim}
python behav_check.py --model Qwen/Qwen2.5-1.5B --n 200 --seed 0

python truth_probe.py signal --dataset counterfact --max-pairs 250 --baseline \
       --perm 200 --rev-counterfact c945b082ca08d0a8f3ba227fb78404a09614c36e
\end{verbatim}
The same two commands with \texttt{--model} set to \texttt{Qwen/Qwen2.5-3B},
\texttt{meta-llama/Llama-3.2-1B}, or \texttt{meta-llama/Llama-3.2-3B}
reproduce the corresponding rows of Table~\ref{b:tab:models}.

\paragraph{Controlled falsification experiments (Table~\ref{tab:seven}, 1--4).}
\begin{verbatim}
# experiment 1 (PHASE ablation) is the PHASE column of `signal` above

# experiment 2: global depth signature (fixed axis), curated pairs
python truth_probe.py domino --dataset builtin --perm 100

# experiment 3: dynamic shade on hard data
python truth_probe.py domino --dataset mix --max-pairs 250 --signal-mag --perm 100 \
       --rev-counterfact c945b082ca08d0a8f3ba227fb78404a09614c36e \
       --rev-truthfulqa 741b8276f2d1982aa3d5b832d3ee81ed3b896490

# experiment 4: per-layer native axes + diagnostic ramp
python truth_probe.py domino --dataset mix --max-pairs 250 --per-layer-axis --perm 100 \
       --rev-counterfact c945b082ca08d0a8f3ba227fb78404a09614c36e \
       --rev-truthfulqa 741b8276f2d1982aa3d5b832d3ee81ed3b896490
\end{verbatim}

\paragraph{Visual observations (Figs.~\ref{fig:coords}--\ref{fig:plane_tqa}, experiments 5--7).}
\begin{verbatim}
python canvas.py --dataset counterfact --layer 16 --out counterfact \
       --rev-counterfact c945b082ca08d0a8f3ba227fb78404a09614c36e
python canvas.py --dataset truthfulqa --layer 16 --out truthfulqa \
       --rev-truthfulqa 741b8276f2d1982aa3d5b832d3ee81ed3b896490
\end{verbatim}
Each \texttt{canvas.py} run writes \texttt{<out>\_coords.png},
\texttt{<out>\_plane.png}, and \texttt{<out>\_3d.png}. Run CounterFact and
TruthfulQA \emph{separately}; do not pass \texttt{--dataset mix} to
\texttt{canvas.py}, as mixing injects a spurious sentence-format axis.

\paragraph{Anatomy (Tables~\ref{tab:peak}--\ref{tab:abl}).}
\begin{verbatim}
# per-layer attention/FFN readability (identity check embedded)
python anatomy.py --dataset counterfact --max-pairs 250 --perm 100 \
       --rev-counterfact c945b082ca08d0a8f3ba227fb78404a09614c36e

# stability of the alternation across 5 seeds
python anatomy_consolidate.py --dataset counterfact --max-pairs 250 --seeds 5 \
       --rev-counterfact c945b082ca08d0a8f3ba227fb78404a09614c36e

# the truth axis inside the expanded FFN space (8960-dim)
python anatomy_expanded.py --dataset counterfact --max-pairs 250 --perm 100 \
       --rev-counterfact c945b082ca08d0a8f3ba227fb78404a09614c36e

# causal ablation, the three bands of Table 5
python ablation.py --dataset counterfact --max-pairs 250 \
       --band-start 11 --band-end 15 --readout 15 \
       --rev-counterfact c945b082ca08d0a8f3ba227fb78404a09614c36e
python ablation.py --dataset counterfact --max-pairs 250 \
       --band-start 16 --band-end 18 --readout 18 \
       --rev-counterfact c945b082ca08d0a8f3ba227fb78404a09614c36e
python ablation.py --dataset counterfact --max-pairs 250 \
       --band-start 19 --band-end 26 --readout 27 \
       --rev-counterfact c945b082ca08d0a8f3ba227fb78404a09614c36e
\end{verbatim}
Anatomy runs use \texttt{float32} (Sec.~\ref{sec:decomp}); about $6$ GB of
VRAM is enough for Qwen2.5-1.5B.

\paragraph{Data provenance and safety.} The curated pairs and the
$2{\times}2$ polarity set are embedded in \texttt{truth\_probe.py}. The
external datasets are \texttt{NeelNanda/counterfact-tracing} (commit
\texttt{c945b08\ldots}) and \texttt{truthful\_qa}, \emph{generation} config
(commit \texttt{741b827\ldots}), loaded as Parquet only. For bit-exact
reproduction the dataset commits may be pinned as shown, or a local file
passed with \texttt{--file-counterfact <path>}. Add \texttt{-h} to any
subcommand for its options.



\paragraph{Norm-weighting control (Methods of Part~I).}
\begin{verbatim}
# repeat with each model and its peak hidden level:
#   Qwen2.5-1.5B->16  Qwen2.5-3B->17  Llama-3.2-1B->8  Llama-3.2-3B->10
python axis_norm_check.py --model Qwen/Qwen2.5-1.5B --layer 16
\end{verbatim}

\section{Reproducibility, Part II}
\label{b:sec:repro}

All tools import \texttt{truth\_probe.py} and \texttt{anatomy.py} from the
same repository and embed the identity, null, and stability
safeguards described above; models load with safetensors only and
\texttt{trust\_remote\_code=False}; the dataset revision is pinned in
Sec.~\ref{b:sec:setup}. Two rules learned inside this project and now part of
its method: the model is passed explicitly to every command (tool defaults
differ, and one default produced the accidental replication of
Sec.~\ref{b:sec:correction}), and no tool prints an automatic verdict. The
Llama models require accepting the license on Hugging Face. Representative
commands; repeat with \texttt{--model} and the landmarks of
Table~\ref{b:tab:models} for each model:

\begin{verbatim}
# behavioral label BEFORE geometry, and the axis/probe gap (Table 1)
python behav_check.py --model meta-llama/Llama-3.2-3B --n 200 --seed 0
python truth_probe.py signal --model meta-llama/Llama-3.2-3B \
       --dataset counterfact --max-pairs 250 --baseline --perm 200 \
       --rev-counterfact c945b082ca08d0a8f3ba227fb78404a09614c36e

# the flip, consolidated: 5 seeds + permutation null + rotation check (Table 2)
python flip_consolidate.py --model meta-llama/Llama-3.2-3B \
       --axis-block 9 --flip-layer 10 --scan-start 5 --scan-end 15

# axis provenance: post and pre frames, both components (Tables 3-4)
python axis_provenance.py --model meta-llama/Llama-3.2-3B --peak 9 \
       --scan-start 5 --scan-end 15
python axis_provenance.py --model meta-llama/Llama-3.2-3B --peak 9 \
       --scan-start 5 --scan-end 15 --component attn

# causal ablation with fixed-frame and frozen controls (Table 5)
python ffn_erosion.py ablate --model meta-llama/Llama-3.2-3B \
       --band-start 10 --band-end 12 --readout 12

# exact gate/value split and causal freezes (Table 6)
python swiglu.py attrib --model meta-llama/Llama-3.2-3B \
       --axis-block 9 --scan-start 6 --scan-end 13
python swiglu.py gatefreeze --model meta-llama/Llama-3.2-3B \
       --band-start 10 --band-end 12 --readout 12

# gauge-invariant circuit norms (weights only, CPU; Table 7)
python circuits.py --model meta-llama/Llama-3.2-3B --dtype bfloat16 \
       --peak 9 --flip 10

# category geometry, arrangement law, dictionary export (Sec. 9-10)
python categories.py --model Qwen/Qwen2.5-3B --peak 16 --write-layer 17 \
--k-relations 8
python categories.py --model meta-llama/Llama-3.2-3B --peak 9 --write-layer 10 --k-relations 8
python arrangement_law.py
python dictionary_export.py --model Qwen/Qwen2.5-3B --peak 16 --write-layer 17 \
--k-relations 8
python dictionary_export.py --model meta-llama/Llama-3.2-3B --peak 9 \
--write-layer 10 --k-relations 8
\end{verbatim}

\texttt{arrangement\_law.py} carries the canonical matrices transcribed in
its source with the regeneration commands in its header; the transcription
was verified cell-by-cell against fresh \texttt{categories.py} runs before
the statistics reported here.

\ Scripts are in \texttt{src/} of the
repository (\texttt{github.com/Francesco-Marhel/TruthProbe}); \
figures
\texttt{flip\_collapse.png}, \texttt{provenance\_asymmetry.png},
\texttt{categories\_geometry.png}, \texttt{arrangement\_law.png} accompany
the source. License: CC BY 4.0.




\begin{thebibliography}{99}

\bibitem{burns2023} C.~Burns, H.~Ye, D.~Klein, J.~Steinhardt.
\emph{Discovering Latent Knowledge in Language Models Without Supervision.}
arXiv:2212.03827, 2022 (ICLR 2023).

\bibitem{marks2024} S.~Marks, M.~Tegmark.
\emph{The Geometry of Truth: Emergent Linear Structure in Large Language
Model Representations of True/False Datasets.}
arXiv:2310.06824, 2023 (COLM 2024).

\bibitem{li2023} K.~Li, O.~Patel, F.~Vi\'egas, H.~Pfister, M.~Wattenberg.
\emph{Inference-Time Intervention: Eliciting Truthful Answers from a
Language Model.} NeurIPS, 2023.

\bibitem{zou2023} A.~Zou et al.
\emph{Representation Engineering: A Top-Down Approach to AI Transparency.}
arXiv:2310.01405, 2023.

\bibitem{burger2024} L.~B\"urger, F.~A.~Hamprecht, B.~Nadler.
\emph{Truth is Universal: Robust Detection of Lies in LLMs.}
arXiv:2407.12831, 2024 (NeurIPS 2024).

\bibitem{cunningham2023} H.~Cunningham, A.~Ewart, L.~Riggs, R.~Huben,
L.~Sharkey. \emph{Sparse Autoencoders Find Highly Interpretable Features in
Language Models.} arXiv:2309.08600, 2023.

\bibitem{bricken2023} T.~Bricken et al.
\emph{Towards Monosemanticity: Decomposing Language Models With Dictionary
Learning.} Transformer Circuits Thread, Anthropic, 2023.

\bibitem{elhage2022} N.~Elhage et al.
\emph{Toy Models of Superposition.}
Transformer Circuits Thread, Anthropic, 2022 (arXiv:2209.10652).

\bibitem{elhage2021} N.~Elhage et al.
\emph{A Mathematical Framework for Transformer Circuits.}
Transformer Circuits Thread, Anthropic, 2021.

\bibitem{barshalom2025} G.~Bar-Shalom, F.~Frasca, Y.~Galron, Y.~Ziser,
H.~Maron. \emph{Beyond Token Probes: Hallucination Detection via Activation
Tensors with ACT-ViT.} arXiv:2510.00296, 2025.

\bibitem{counterfact} K.~Meng, D.~Bau, A.~Andonian, Y.~Belinkov.
\emph{Locating and Editing Factual Associations in GPT.}
NeurIPS, 2022. Dataset mirror: \texttt{NeelNanda/counterfact-tracing}.

\bibitem{meng2022} K.~Meng, D.~Bau, A.~Andonian, Y.~Belinkov.
\emph{Locating and Editing Factual Associations in GPT.}
NeurIPS, 2022. Dataset mirror: \texttt{NeelNanda/counterfact-tracing}.

\bibitem{icr2025} \emph{ICR Probe: Tracking Hidden State Dynamics for
Reliable Hallucination Detection in LLMs.} arXiv:2507.16488, 2025.

\bibitem{hsad2025} \emph{LLM Hallucination Detection: A Fast Fourier
Transform Method Based on Hidden Layer Temporal Signals (HSAD).}
arXiv:2509.13154, 2025.

\bibitem{rotational2026} \emph{How Transformers Reject Wrong Answers:
Rotational Dynamics of Factual Constraint Processing.}
Preprint, February 2026.

\bibitem{headgenome} \emph{The HeadGenome Taxonomy: A Structural and Causal
Map of Attention Head Specialization in Transformers.}
Preprint, OpenReview iG4i2pc4wg.

\bibitem{platonic2024} M.~Huh, B.~Cheung, T.~Wang, P.~Isola.
\emph{The Platonic Representation Hypothesis.}
ICML 2024. arXiv:2405.07987.

\bibitem{truthfulqa} S.~Lin, J.~Hilton, O.~Evans.
\emph{TruthfulQA: Measuring How Models Mimic Human Falsehoods.} ACL, 2022.

\end{thebibliography}

\end{document}
